\documentclass[11pt]{article}
\usepackage{amsmath,amssymb,amsthm,mathtools}
\usepackage[margin=0.85in]{geometry}
\usepackage{enumitem}
\usepackage[colorlinks=true,linkcolor=blue,citecolor=blue]{hyperref}

\newtheorem{theorem}{Theorem}[section]
\newtheorem{lemma}[theorem]{Lemma}
\newtheorem{proposition}[theorem]{Proposition}
\newtheorem{corollary}[theorem]{Corollary}
\theoremstyle{remark}
\newtheorem{remark}[theorem]{Remark}

\newcommand{\R}{\mathbb{R}}
\newcommand{\E}{\mathbb{E}}
\newcommand{\F}{\mathcal{F}}
\newcommand{\W}{\mathbb{W}}
\newcommand{\one}{\mathbf{1}}
\newcommand{\norm}[1]{\lVert #1\rVert}
\newcommand{\abs}[1]{\lvert #1\rvert}
\newcommand{\checkK}[1]{\check{#1}}
\newcommand{\inner}[2]{\langle #1,#2\rangle}
\newcommand{\SET}[1]{\{#1\}}
\DeclareMathOperator{\diver}{div}
\DeclareMathOperator{\supp}{supp}

\begin{document}

\section{Problem and answer}

Let $D\subseteq\R^{d}$, $d\ge2$, be the Euclidean unit ball. For a
H\"older--Zygmund conductivity $f\in C^{\alpha}(D)$, $\alpha>1$, let
$u=u_{f}$ solve the Darcy problem
\[
\diver(f\nabla u)=1\ \ \text{on }D,\qquad u|_{\partial D}=0 .
\]
We observe
\[
Y\;=\;u_{f}\;+\;\frac{1}{\sqrt{n}}\,\W ,
\]
where $\W$ is Gaussian white noise on $L^{2}(D)$, and we estimate $f$ over
the class
\[
\F(\alpha,r)=\SET{\,h\in C^{\alpha}(D):\ \norm{h-\one}_{C^{\alpha}(D)}\le r\,},
\qquad r\in(0,1)\ \text{small and fixed}.
\]
(The class is contained in $C^{a}$ with the $a$-th derivatives uniformly
$t$-H\"older, where $\alpha=a+t$, $a=\lceil\alpha\rceil-1\in\mathbb Z_{\ge1}$,
$t\in(0,1]$, Zygmund if $t=1$; in particular every $f\in\F(\alpha,r)$ is
continuously differentiable and satisfies $1-r\le f\le1+r$.)

\begin{theorem}\label{thm:main}
There exist $r_{0}=r_{0}(\alpha,d)\in(0,1]$ and, for every fixed
$r\in(0,r_{0}]$, constants $c,C>0$ with $C$ depending only on
$(\alpha,d)$ and $c$ on $(\alpha,d,r)$, such that
\[
c\,\Phi_{n}
\ \le\
\inf_{\widehat f_{n}}\ \sup_{f_{0}\in\F(\alpha,r)}
\E_{f_{0}}\!\bigl[\norm{\widehat f_{n}-f_{0}}_{L^{2}(D)}^{2}\bigr]
\ \le\ C\,\Phi_{n},
\qquad
\Phi_{n}=n^{-\frac{2\alpha}{2\alpha+d+2}} .
\]
\end{theorem}

\noindent\textbf{Answer.} The minimax rate is
\[
\boxed{\ n^{-\frac{2\alpha}{2\alpha+d+2}}\ }
\qquad\text{up to logarithmic factors.}
\]
Both bounds proved below carry this power of $n$; our argument is in fact
log-free, which is a fortiori ``up to logarithmic factors.'' This is
\emph{faster} than the
direct white-noise rate $n^{-2\alpha/(2\alpha+d)}$: the denominator gains
an additive $+2$, i.e.\ the inverse problem is only \emph{mildly ill-posed
of degree one}. The mechanism, proved below, is that the linearization $T$
of $f\mapsto u_{f}$ at $\one$ satisfies $\Delta(Tg)=g+\frac{x\cdot\nabla}{d}g$,
and the transport-type source $g\mapsto g+\frac{x\cdot\nabla}{d}g$ is
invertible as an order-zero operator ($L^{-1}=J\in\mathcal B(L^{2}(D))$ with
$\norm{J}\le2d$, $L\in\mathcal B(H^{1}(D),L^{2}(D))$;
Lemma~\ref{lem:Linv}); consequently $T$
behaves like a general (nonlocal) operator of order $-1$, gaining exactly
one derivative on scale-$\delta$ perturbations, rather than two.

\noindent The lower bound is Proposition~\ref{prop:lower}; the upper bound
is Proposition~\ref{prop:upper}. Constants implicit in $\lesssim,\asymp$ depend only on
$(\alpha,d)$ and a fixed bump, never on $f$, $n$, or $\delta$; $C$ is
uniform in $r\in(0,r_0]$, while the lower-bound $c$ carries the factor
$r^{2}$ of the shrinking class (Proposition~\ref{prop:lower}).

\section{Preliminaries}\label{sec:prelim}

\subsection{White-noise model and likelihood}
Gaussian white noise $\W$ on $L^{2}(D)$ is the centred Gaussian process
indexed by $L^{2}(D)$ with
$\E\inner{\W}{\phi}=0$ and
$\E[\inner{\W}{\phi}\inner{\W}{\psi}]=\inner{\phi}{\psi}_{L^{2}(D)}$. Under
conductivity $f$, the law $P_{f}$ of $Y$ is the Gaussian shift with mean
$u_{f}$ and covariance $n^{-1}I$ on $L^{2}(D)$. For two conductivities
$f,f_0$ the Cameron--Martin formula gives
\begin{equation}\label{eq:KL}
\operatorname{KL}(P_{f}\,\|\,P_{f_{0}})
\;=\;\frac{n}{2}\,\norm{u_{f}-u_{f_{0}}}_{L^{2}(D)}^{2}.
\end{equation}
All norms $\norm{\cdot}$ are $L^{2}(D)$ unless otherwise indicated;
$\norm{\cdot}_{p}$ denotes $L^{p}(D)$.

\subsection{External tools (precise statements)}
We collect the standard results used, with precise hypotheses, so that no
nonstandard result enters tacitly.

\begin{enumerate}[label=(T\arabic*),itemsep=1pt,leftmargin=2.2em]
\item\label{T:fano}\textbf{Fano's inequality.} For probability measures
$P_1,\dots,P_N$ on a common space, any estimator $\widehat\theta$ of the
index, and any reference $\bar P$,
$\frac1N\sum_{j} P_j(\widehat\theta\ne j)\ge 1-\frac{\frac1N\sum_j
\operatorname{KL}(P_j\|\bar P)+\log 2}{\log N}$.
\item\label{T:VG}\textbf{Varshamov--Gilbert.} For $M\ge1$ there exist
$N\ge 2^{cM}$ subsets $S\subset\SET{1,\dots,M}$, each of cardinality $M/2$,
with pairwise Hamming distance $\abs{S\triangle S'}\ge M/8$; $c>0$
absolute.
\item\label{T:weyl}\textbf{Weyl's law.} For the Dirichlet Laplacian on the
bounded smooth domain $D$, $\#\SET{\lambda_\nu\le\Lambda}\sim
c_d\abs D\Lambda^{d/2}$; hence the span of the eigenfunctions with
$\lambda_\nu\le M^2$ has dimension $\asymp M^{d}$.
\item\label{T:CM}\textbf{Cameron--Martin.} For $Y=u+n^{-1/2}\W$ on
$L^2(D)$, means $u_1,u_0\in L^2(D)$ give laws with
$\operatorname{KL}(P_{u_1}\|P_{u_0})=\frac{n}{2}\norm{u_1-u_0}_2^2$.
\item\label{T:wave}\textbf{Wavelet/Besov characterization of $C^\alpha$
and square function on $D$.} There is a boundary-corrected wavelet frame of
the smooth domain $D$ (Cohen--Dahmen--DeVore-type) such that for
$\alpha>0$ (noninteger, with the Zygmund convention at integers),
$f\in C^\alpha(D)=B^\alpha_{\infty,\infty}(D)$ iff
$\sup_{\lambda}2^{j_\lambda(\alpha+d/2)}\abs{\inner{f}{\psi_\lambda}}
<\infty$, comparable to $\norm{f}_{C^\alpha}$; and
$\norm{f}_{L^2}\asymp\bigl(\sum_j\norm{\Delta_j f}_{L^2}^2\bigr)^{1/2}$ for
the spectral dyadic pieces $\Delta_j=\varphi_j(\sqrt{-\Delta})$ (the
Dirichlet-Laplacian smooth functional calculus / heat-kernel Mikhlin
square-function inequality on the smooth bounded domain).
\item\label{T:schauder}\textbf{Divergence-form elliptic regularity.} For
$1-r\le f\le1+r$, $f\in C^\alpha$, $\alpha>1$, $\diver(f\nabla u)=F$ in
$D$ with $u=0$ on $\partial D$ has a unique solution with
$\norm{u}_{C^{1,\beta}}\le C(\alpha,d,r)\norm{F}_{L^\infty}$ for some
$\beta=\beta(\alpha,d)\in(0,1]$; in particular
$\norm{\nabla u}_{L^\infty}\le C$.
\item\label{T:hardy}\textbf{Hardy / Ces\`aro.} The Ces\`aro operator
$(Cf)(t)=\frac1t\int_0^t f(s)\,ds$ has norm $\frac{p}{p-1}$ ($=2$ on $L^2$)
on $L^p((0,1),dt)$.
\end{enumerate}

\noindent Every PDE/functional bound below is either derived here or is an
application of (T1)--(T7).

\subsection{Linearization about the constant conductivity}
The constant-conductivity solution is explicit:
\[
u_{0}(x)=\frac{\abs{x}^{2}-1}{2d},
\qquad \nabla u_{0}(x)=\frac{x}{d},
\qquad \Delta u_{0}=1 .
\]
For any $\phi\in C^{1}(D)$ one has the pointwise identity
\begin{equation}\label{eq:u0prop}
\diver(\phi\nabla u_{0})
=\nabla\phi\cdot\Bigl(\frac{x}{d}\Bigr)+\phi\,\Delta u_{0}
=\frac{x\cdot\nabla\phi}{d}+\phi .
\end{equation}

\begin{lemma}[Linearization]\label{lem:lin}
For $g\in L^{2}(D)$, let $Tg\in H^{1}_{0}(D)$ be the unique solution of
$-\diver(\nabla w)=-\diver(g\nabla u_{0})$, $w|_{\partial D}=0$, i.e.\
$Tg=-\Delta^{-1}\!\bigl(\diver(g\nabla u_{0})\bigr)$ in the
$H^{-1}(D)$/$H^{1}_{0}(D)$ sense. Then, as distributions,
\begin{equation}\label{eq:DeltaT}
\Delta(Tg)=g+\frac{x\cdot\nabla}{d}\,g .
\end{equation}
\end{lemma}
\begin{proof}
By \eqref{eq:u0prop}, $\diver(g\nabla u_{0})=g+\frac{x\cdot\nabla}{d}g$
as a distribution of order $\le1$, hence in $H^{-1}(D)$; applying
$\Delta^{-1}$ gives \eqref{eq:DeltaT}. The Dirichlet Laplacian
$-\Delta:H^{1}_{0}(D)\to H^{-1}(D)$ is an isomorphism (Lax--Milgram), and
$\norm{Tg}_{H^{1}}\lesssim\norm{g}_{H^{-1}}$.
\end{proof}

\begin{remark}[Transport structure of $L$]\label{rem:meanzero}
Write $A=\frac{x\cdot\nabla}{d}$ and $L=I+A$. Using
$\diver(xg)=(\diver x)g+x\cdot\nabla g=dg+x\cdot\nabla g$ and
$x\cdot\nu=\abs{x}=1$ on $\partial D$,
\[
Lg=g+\frac{x\cdot\nabla g}{d}
=\frac1d\bigl(dg+x\cdot\nabla g\bigr)
=\frac1d\diver(x\,g).
\]
Hence $L$ is a pure divergence composed with the vector field $x\mapsto
x/d$; in particular \emph{if} $\int_D g\,dx=0$ \emph{then}
$\int_D Lg\,dx=\frac1d\int_{\partial D}g\,dS$ (this vanishes for the
mean-zero wavelet packets used below since each $\psi_m$ has integral zero
and is interior). This divergence/transport structure makes $L$ an
order-zero invertible operator on
$L^{2}(D)$ (Lemma~\ref{lem:Linv}); the smoothing order of $T=-\Delta^{-1}L$
is then governed by the dilation generator, giving exactly one degree
(Lemmas~\ref{lem:fourier}, \ref{lem:scale2}).
\end{remark}

\subsection{The transport operator $L$ is an invertible order-zero operator}

\begin{lemma}[$L$ is invertible, $\norm{L^{-1}}\le C_{d}$]\label{lem:Linv}
Let $A=\frac{x\cdot\nabla}{d}$ and $L=I+A$. For $g\in C^{1}(D)$ define
\[
(Jg)(x)\;:=\;d\int_{0}^{1}\tau^{d-1}g(\tau x)\,d\tau ,
\qquad x\in D\setminus\SET{0}.
\]
Then:
\begin{enumerate}[label=(\roman*),itemsep=1pt]
\item $LJg=JLg=g$ on $C^{1}(D)$, so $L$ is invertible and $L^{-1}=J$.
\item $\norm{L^{-1}g}_{L^{2}(D)}\le C_{d}\norm{g}_{L^{2}(D)}$ for every
$d\ge2$ (in fact $C_d=2d$ suffices).
\item Consequently $\norm{Lg}_{2}\ge c_d\norm{g}_{2}$ ($c_d=C_d^{-1}>0$).
Precisely: $J=L^{-1}$ extends to an element of $\mathcal B(L^{2}(D))$,
$L\in\mathcal B(H^{1}(D),L^{2}(D))$ (since
$\norm{Lg}_2\le\norm{g}_2+\frac1d\norm{\nabla g}_2$). Moreover, writing
$Lg=\frac1d\diver(xg)$ and defining $\diver(xg)\in H^{-1}(D)$ for general
$g\in L^{2}(D)$ by duality, the operator
$T(g):=-\Delta^{-1}\bigl(\tfrac1d\diver(xg)\bigr)$ extends the
$\Delta w=Lg$ map of Lemma~\ref{lem:lin} from $C^{1}$ to all of
$L^{2}(D)$, with
\[
T\in\mathcal B\bigl(L^{2}(D),H^{1}_{0}(D)\bigr)
\cap\mathcal B\bigl(H^{-1}(D),L^{2}(D)\bigr),
\qquad\norm{T}_{L^{2}\to L^{2}}\le C_{d}\cdot 2C_{P},
\]
in particular $\norm{Tg}_{2}\le C_{T}\norm{g}_{2}$ with
$C_T=C_{T}(d)$, used in the spectral-tail step of
Lemma~\ref{lem:invert}.
\end{enumerate}
\end{lemma}
\begin{proof}
(i) Fix $x\ne0$ and set $\phi(\tau)=g(\tau x)$, so
$\phi'(\tau)=x\cdot\nabla g(\tau x)$. Then
\[
\frac{d}{d\tau}\bigl[\tau^{d}g(\tau x)\bigr]
=d\tau^{d-1}g(\tau x)+\tau^{d}x\cdot\nabla g(\tau x) .
\]
Integrating over $\tau\in[0,1]$ ($\tau^{d}g(\tau x)\to0$ as $\tau\to0$):
\[
g(x)=\int_{0}^{1}\Bigl(d\tau^{d-1}g(\tau x)+\tau^{d}x\cdot\nabla g(\tau x)\Bigr)d\tau
\]
and, differentiating under the integral in $x$,
\[
x\cdot\nabla(Jg)(x)
=d\int_0^1\tau^{d-1}\bigl(\tau x\cdot\nabla g(\tau x)\bigr)d\tau
=d\int_0^1\tau^{d}\,x\cdot\nabla g(\tau x)\,d\tau .
\]
Adding,
$d\,(Jg)(x)+x\cdot\nabla(Jg)(x)
=d\int_0^1\bigl(d\tau^{d-1}g(\tau x)+\tau^{d}x\cdot\nabla g(\tau x)\bigr)d\tau
=d\,g(x)$ by the previous display; dividing by $d$,
$L(Jg)(x)=(Jg)(x)+\frac{x\cdot\nabla(Jg)(x)}{d}=g(x)$, i.e.\ $LJ=I$.
Applying $L$ first and using the same integration with
$(Lg)(\tau x)=g(\tau x)+\frac{\tau x\cdot\nabla g(\tau x)}{d}$ gives
$JLg=g$; hence $L^{-1}=J$.

(ii) Polar coordinates $x=\rho\theta$, $dx=\rho^{d-1}d\rho\,d\sigma(\theta)$:
\[
\norm{L^{-1}g}_{2}^{2}
=\int_{\mathbb S^{d-1}}\int_{0}^{1}
\Bigl|d\int_{0}^{1}\tau^{d-1}g(\tau\rho\theta)d\tau\Bigr|^{2}
\rho^{d-1}d\rho\,d\sigma(\theta).
\]
Substitute $s=\tau\rho$:
$d\int_0^1\tau^{d-1}g(\tau\rho\theta)d\tau
=\frac{d}{\rho^{d}}\int_0^{\rho}s^{d-1}g(s\theta)\,ds$. For fixed
$\theta\in\mathbb S^{d-1}$ set $\gamma(\rho)=g(\rho\theta)$ and equip
$(0,1)$ with the measure $d\mu(\rho)=\rho^{d-1}d\rho$, so that
$\mu(0,\rho)=\rho^{d}/d$ and
$\norm{g}_{L^{2}(D)}^{2}=\int_{\mathbb S^{d-1}}\norm{\gamma}_{L^{2}(d\mu)}^{2}
d\sigma(\theta)$. The pointwise formula reads
\[
(L^{-1}g)(\rho\theta)
=\frac{d}{\rho^{d}}\int_0^{\rho}\gamma(s)\,d\mu(s)
=d\,(H\gamma)(\rho),
\qquad
(H\gamma)(\rho):=\frac{1}{\mu(0,\rho)}\int_0^{\rho}\gamma\,d\mu .
\]
The averaging operator $H$ is bounded on $L^{2}(d\mu)$
(Hardy's inequality on $((0,1),d\mu)$). To compute a constant, change
variables $t=\rho^{d}$: then $\rho^{d-1}d\rho=dt/d$, $\mu(0,\rho)=t/d$, and
with $\beta(t):=d^{-1/2}\gamma(t^{1/d})$ one has the isometry
$\int_0^1|\beta(t)|^{2}dt=\int_0^1|\gamma(\rho)|^{2}\rho^{d-1}d\rho
=\norm{\gamma}_{L^{2}(d\mu)}^{2}$, and
\[
(H\gamma)(\rho)=\frac{d}{t}\int_{0}^{t^{1/d}}\gamma(s)\,s^{d-1}ds
=\frac{d}{t}\int_{0}^{t}\gamma(s_1^{1/d})\,\frac{ds_1}{d}
=\frac{\sqrt d}{t}\int_{0}^{t}\beta(s_1)\,ds_1 ,
\]
which is $\sqrt d$ times the Ces\`aro operator in $t$ (of norm
$\frac{p}{p-1}=2$ on $L^{2}((0,1),dt)$). Therefore
\[
\norm{H\gamma}_{L^{2}(d\mu)}^{2}
=\int_0^1\abs{H\gamma}^{2}\frac{dt}{d}
=\frac1d\cdot d\int_0^1\Bigl|\frac1t\int_0^t\beta\Bigr|^{2}dt
\le\frac{d}{d}\cdot4\int_0^1\abs{\beta}^{2}dt
=4\,\norm{\gamma}_{L^{2}(d\mu)}^{2},
\]
so $\norm{H\gamma}_{L^{2}(d\mu)}\le2\norm{\gamma}_{L^{2}(d\mu)}$.
Therefore
\[
\norm{L^{-1}g}_{L^{2}(D)}^{2}
=d^{2}\!\!\int_{\mathbb S^{d-1}}\!\!\norm{H\gamma}_{L^{2}(d\mu)}^{2}d\sigma
\le d^{2}\cdot 4\int_{\mathbb S^{d-1}}\norm{\gamma}_{L^{2}(d\mu)}^{2}d\sigma
=4d^{2}\,\norm{g}_{L^{2}(D)}^{2},
\qquad
\norm{L^{-1}g}_{2}\le2d\,\norm{g}_{2}.
\]
Thus $\norm{L^{-1}}\le C_{d}$ with $C_d=2d$; any $d$-dependent
constant suffices here, as all constants below are allowed to depend on
$(\alpha,d)$. We thus take (ii) in the form
$\norm{L^{-1}g}_{2}\le C_d\norm{g}_{2}$.

(iii) From $\norm{g}_{2}=\norm{L^{-1}Lg}_{2}\le C_d\norm{Lg}_{2}$,
$\norm{Lg}_{2}\ge c_d\norm{g}_{2}$. On a dense set ($C^{1}$) $J=L^{-1}$
with $\norm J\le C_d$, so $J$ extends to $\mathcal B(L^{2}(D))$; and
$\norm{Lg}_{2}\le\norm{g}_{2}+\frac1d\norm{x\cdot\nabla g}_{2}
\le(1+\frac1d)\norm{g}_{H^{1}}$, so $L\in\mathcal B(H^{1},L^{2})$ (it is
\emph{not} bounded on $L^{2}$). For the last claim, map
$g\in L^{2}$ to the divergence-form distribution $F=\frac1d\diver(xg)\in
H^{-1}$ ($\norm{F}_{H^{-1}}\le\frac1d\norm{g}_2$) and apply the Dirichlet
Laplacian inverse: $-\Delta^{-1}:H^{-1}\to H^{1}_{0}$ is bounded, so
$T:L^{2}\to H^{1}_{0}$ is bounded with
$\norm{Tg}_{H^{1}}\lesssim\frac1d\norm{g}_2$; by Poincar\'e
$\norm{Tg}_{2}\lesssim\norm{g}_{2}$. Finally, for $\psi\in L^{2}$ and
$\phi=-\Delta^{-1}\psi\in H^{2}\cap H^{1}_{0}$
($\norm{\phi}_{H^{2}}\lesssim\norm{\psi}_2$),
$\abs{\inner{Tg}{\psi}}=\frac1d\abs{\int_D g\,x\cdot\nabla\phi\,dx}
\le\frac1d\norm{g}_{H^{-1}}\norm{x\cdot\nabla\phi}_{H^{1}_0}
\lesssim\norm{g}_{H^{-1}}\norm{\psi}_2$, so
$T:H^{-1}\to L^{2}$ is bounded.
On $g\in C^{1}$ the identity $\frac1d\diver(xg)=Lg$ identifies this
extension with the $T$ of Lemma~\ref{lem:lin}; in particular $Tg$ always
satisfies the very-weak form of $\Delta w=Lg$. The refined scale-localized
bound $\norm{Tg}_{2}\lesssim\delta\norm{g}_{2}$ for packets is
Lemma~\ref{lem:scale2}.
\end{proof}

\begin{remark}\label{rem:crux}
Lemma~\ref{lem:Linv} is the analytic crux: by Lemma~\ref{lem:lin},
$Tg=-\Delta^{-1}Lg$. On a frequency band $\abs{\cdot}\asymp\delta^{-1}$ the
Dirichlet Laplacian inverts with gain $\delta^{2}$ but the \emph{transport
source} $L$ is invertible as an order-zero operator (Lemma~\ref{lem:Linv}:
$L^{-1}=J\in\mathcal B(L^{2})$ and $\norm{Lg}_2\ge c_d\norm{g}_2$), so $T$ gains only
\emph{one} power of $\delta$. This is made precise in
Lemma~\ref{lem:scale2}.
\end{remark}

\begin{lemma}[Elliptic estimate for the harmonic
part]\label{lem:harmonic}
Let $F\in H^{-1}(D)$ and decompose the solution $w\in H^{1}_{0}(D)$ of
$-\Delta w=F$, $w|_{\partial D}=0$, as $w=U-q$ with
$U=(-\Delta_{\R^{d}})^{-1}F\in H^{1}_{\rm loc}$ (fundamental-solution
representation on the compactly supported $F$) and $\Delta q=0$ in $D$,
$q=U$ on $\partial D$. Then
\[
\norm{q}_{L^{2}(D)}\ \le\ C_{E}\,\norm{U}_{H^{1}(D)},
\qquad C_E=C_E(d),
\]
and consequently, for $F=\frac1d\diver(xg)$ ($g\in L^{2}$) and
$w=Tg$,
$\norm{q}_{2}\le C_E''\norm{g}_{2}$ with $C_E''=C_E''(d)$.

\noindent If moreover $\supp g\subset D_{(\tau)}:=\SET{\abs x\le1-\tau}$
for some $\tau\in(0,\tfrac12]$, then with $C_E'''=C_E'''(d)$,
\[
\norm{U}_{H^{1}(\partial D)}\ \le\ C_E'''\,\tau^{-(d-1)}\norm{g}_2 ,
\]
i.e.\ the harmonic trace stays $O(\norm g_2)$ (up to a
$\tau$-\emph{fixed} factor) and is not amplified as the source recedes from
$\partial D$.
\end{lemma}
\begin{proof}
Set $\phi:=q-U$, the solution of $-\Delta\phi=F$ in $H^{1}_{0}(D)$
(since $\Delta q=0$, $-\Delta U=F$, and $\phi=q-U=0$ on $\partial D$);
$H^{1}_0$ energy theory (Lax--Milgram and Poincar\'e) gives
$\norm{\phi}_{H^{1}}\le C_P\norm{F}_{H^{-1}}$. Testing against $\phi$,
\[
\int_D\abs{\nabla\phi}^{2}
=\inner{F}{\phi}_{H^{-1},H^{1}_0}
\le\norm{F}_{H^{-1}}\norm{\phi}_{H^{1}},
\]
so $\norm{\phi}_{H^{1}}\le C_P\norm{F}_{H^{-1}}$, and
$\norm{q}_{L^{2}}\le\norm{U}_{L^{2}}+\norm{\phi}_{L^{2}}
\le\norm{U}_{L^{2}}+C_P\norm{F}_{H^{-1}}
\le C_E\norm{U}_{H^{1}}$ after using
$\norm{F}_{H^{-1}}\lesssim\norm{U}_{H^{1}}$.
For the specific $F=\frac1d\diver(xg)$,
$\norm{U}_{H^{1}(D)}\lesssim\norm{F}_{H^{-1}}\le\frac1d\norm g_2$
(interior elliptic regularity on $D\Subset\tfrac32 D$ for
$-\Delta_{\R^{d}}U=F$, plus $\norm{F}_{H^{-1}}\le\frac1d\norm g_2$ from
the divergence form), so $\norm{q}_{2}\le C_E''\norm{g}_{2}$.
\smallskip
For the interior claim, use the representation
$U=(-\Delta_{\R^d})^{-1}\bigl(\tfrac1d\diver(xg)\bigr)
=-\frac1d(-\Delta_{\R^d})^{-1}\diver(xg)$, i.e.\
$U(x)=-\frac1d\int\nabla\Phi(x-y)\cdot y\,g(y)\,dy$ with $\Phi$ the
fundamental solution, $\abs{\nabla\Phi(z)}\lesssim\abs z^{1-d}$
(integration by parts is unnecessary: this is the convolution of the
$L^{2}$ field $yg$ with the locally-integrable kernel
$\nabla\Phi$). For $x\in D\setminus D_{(\tau/2)}$ and
$y\in D_{(\tau)}$, $\abs{x-y}\ge\tau/2$, so
$\abs{\nabla\Phi(x-y)}\lesssim\tau^{1-d}$ and
$\abs{U(x)}\lesssim\tau^{1-d}\norm{g}_1\lesssim\tau^{1-d}\norm g_2$,
and the same with one more derivative,
$\abs{\nabla U(x)}\lesssim\tau^{-d}\norm g_2$. Hence for a fixed
interior/annular collar $D\setminus D_{(\tau/2)}$,
$\norm{U}_{H^{1}(D\setminus D_{(\tau/2)})}\lesssim\tau^{1-d}\norm g_2$
and by smooth continuation (interior elliptic regularity on
$D_{(\tau/2)}$ for the same equation, applied to the patched field)
$\norm{U}_{H^{1}(D)}\lesssim\tau^{1-d}\norm g_2$; the trace theorem
$\norm{u}_{L^{2}(\partial D)}\lesssim\norm u_{H^{1}(D)}$ then gives
$\norm U_{L^{2}(\partial D)}\lesssim\tau^{1-d}\norm g_2$ with $\nabla U$
similarly bounded, so
$\norm U_{H^{1}(\partial D)}\lesssim\tau^{1-d}\norm g_2$. The stated
exponent $\tau^{-(d-1)}$ is not optimal but suffices: any
$(\tau,d)$-dependent constant works, since $\tau$ is chosen before
$\delta_0$.
\end{proof}

\subsection{The forward map: very-weak form, apriori bounds, quadratic
remainder}\label{sec:forward}

Since $f\in C^{\alpha}$, $\alpha>1$, $f$ is $C^{1}$ and
$1-r\le f\le1+r$ uniformly. We use the very-weak ($L^{2}$) formulation:
$u_f\in L^{2}(D)$ satisfies, for all $\phi\in H^{2}(D)\cap H^{1}_{0}(D)$,
\begin{equation}\label{eq:veryweak}
\int_D u_f\,\diver(f\nabla\phi)\,dx=\int_D\phi\,dx .
\end{equation}
Existence/uniqueness of a very-weak solution follows by transposition:
since $f\in C^{1}$ and $1-r\le f\le1+r$, the operator
$-\diver(f\nabla\cdot):H^{1}_0(D)\to H^{-1}(D)$ is an isomorphism
(Lax--Milgram), so for $\psi\in L^{2}(D)$ the adjoint problem
$-\diver(f\nabla\phi)=\psi$, $\phi\in H^{1}_0(D)$, has a unique solution
$\phi\in H^{2}(D)\cap H^{1}_0(D)$; defining $u_f\in L^{2}(D)$ by
$\langle u_f,\psi\rangle:=\int_D\phi\,dx$ (well-defined, linear, bounded in
$\psi$ since $\norm{\phi}_{H^{2}}\lesssim\norm{\psi}_2$) gives
\eqref{eq:veryweak} upon $\diver(f\nabla\phi)=-\psi$, and uniqueness in
$L^{2}$ follows from the same isomorphy (if $u\in L^{2}$ satisfies
\eqref{eq:veryweak} with the right side $0$, the density of the adjoint
solutions $\psi$ in $L^{2}$ forces $u=0$). All subsequent $L^{2}$/weak
estimates use only ellipticity and boundedness via (T6); we never invoke
$u_f\in H^{1}_0(D)$ from \eqref{eq:veryweak}.

\begin{lemma}[Apriori estimates]\label{lem:apriori}
Uniform in $f\in\F(\alpha,r)$:
(i) $\norm{u_f}_{\infty}+\norm{u_f}_{2}\lesssim1$;
(ii) $\norm{u_f-u_0}_{2}\lesssim\norm{f-\one}_{2}\le r$;
(iii) more generally $\norm{u_f-u_0}_{2}\le C_A\norm{(f-\one)\nabla
u_0}_{H^{-1}(D)}$ with $C_A=C_A(d,r)$, and in particular
$\norm{u_f-u_0}_2\lesssim\norm{f-\one}_\infty$; if moreover
$\supp(f-\one)\subset D'$ with $\operatorname{dist}(D',\partial
D)>\varepsilon_0$, then
$\norm{u_f-u_0}_2\lesssim_{\varepsilon_0}\norm{f-\one}_{H^{1}(D')}$.
\end{lemma}
\begin{proof}
(i) By the comparison principle with barriers
$\frac{u_0}{1\mp r}=\frac{\abs{x}^{2}-1}{2d(1\mp r)}$: since
$1-r\le f\le1+r$ and
$\diver(f\nabla u_f)=1$, the weak maximum principle for the divergence
operator $w\mapsto\diver(f\nabla w)$ (with the subsolution/supersolution
comparison valid in the weak sense) gives, comparing $u_f$ against the
two radial barriers,
$\diver\bigl(f\nabla\tfrac{u_0}{1+r}\bigr)
=\tfrac{f}{1+r}\Delta u_0=\tfrac{f}{1+r}\le1=\diver(f\nabla u_f)$
with $\tfrac{u_0}{1+r}=0=u_f$ on $\partial D$, so
$u_f\ge\tfrac{u_0}{1+r}$ in $D$; symmetrically
$\diver\bigl(f\nabla\tfrac{u_0}{1-r}\bigr)
=\tfrac{f}{1-r}\ge1$ gives $u_f\le\tfrac{u_0}{1-r}$. Hence
$\frac{\abs{x}^{2}-1}{2d(1+r)}\le u_f\le\frac{\abs{x}^{2}-1}{2d(1-r)}\le0$
and
$\abs{u_f}\le\frac{1}{2d(1-r)}\lesssim1$ on $D$.
(ii) Testing \eqref{eq:veryweak} for $u_f$ and $u_0$ against
$\phi=-\Delta^{-1}\psi$ ($\psi\in L^{2}$ arbitrary,
$\norm{\phi}_{H^{2}}\lesssim\norm{\psi}_{2}$) and subtracting:
\[
\inner{u_f-u_0}{\psi}
=-\int_D(f-\one)\nabla u_0\cdot\nabla\phi\,dx ,
\]
so $\abs{\inner{u_f-u_0}{\psi}}\le\norm{f-\one}_{2}\norm{\nabla u_0}_{\infty}
\norm{\nabla\phi}_{2}\lesssim\norm{f-\one}_{2}\norm{\psi}_{2}$, and (ii)
follows by duality.
(iii) Sharpening the same display,
$\abs{\inner{u_f-u_0}{\psi}}
\le\norm{(f-\one)\nabla u_0}_{H^{-1}}\norm{\nabla\phi}_{H^{1}_0}
\le C\norm{(f-\one)\nabla u_0}_{H^{-1}}\norm{\psi}_2$, giving
$\norm{u_f-u_0}_2\lesssim\norm{(f-\one)\nabla u_0}_{H^{-1}}$.
Now $(f-\one)\nabla u_0=(f-\one)x/d$. On the full ball
$\norm{(f-\one)x/d}_{H^{-1}}\le\norm{(f-\one)x/d}_{2}
\lesssim\norm{f-\one}_\infty$ (since $H^{-1}$ is weaker than $L^{2}$ via
the embedding of the source as a divergence).
If $\supp(f-\one)\subset D'\Subset D$, fix a cut-off
$\chi_0\in C^{\infty}_c(D)$, $\chi_0=1$ on $D'$, and compute
$\diver\bigl(\chi_0(f-\one)x/d\bigr)
=(f-\one)\frac{x\cdot\nabla\chi_0}{d}+\chi_0\bigl(f-\one+\frac{x\cdot
\nabla f}{d}\bigr)$; this is an $H^{-1}$ function of norm
$\lesssim\norm{f-\one}_{L^{2}(D')}+\norm{\nabla f}_{L^{2}(D')}$, and
$\norm{(f-\one)x/d}_{H^{-1}}=\norm{\chi_0(f-\one)x/d}_{H^{-1}}
\lesssim\norm{f-\one}_{H^{1}(D')}$ by Lax--Milgram duality
(the map $g\mapsto\diver(\chi_0gx/d)$ is bounded
$H^{1}_{0}(D)\to H^{-1}(D)$ and for $g$ supported in $D'$, $\chi_0g=g$).
(The gradient term $\nabla f$ is present because the conductivity does
\emph{not} vanish on $\partial D$; without the interior-support
hypothesis the map $f\mapsto u_f$ is $C^\alpha\mapsto L^{2}$ bounded only
with the additive $C(f-\one)$ constant of (ii), which is consistent with
-- and does not improve -- the radius minimax rate.)
\end{proof}

\begin{lemma}[Quadratic remainder]\label{lem:quad}
For $f=\one+h\in\F(\alpha,r)$, write
$u_f-u_0=Th+R(h)$. Then
$\norm{R(h)}_{2}\le C\norm{h}_{\infty}\norm{h}_{2}$.
\end{lemma}
\begin{proof}
From the equations for $u_f$ and $u_0$, $w:=u_f-u_0$ satisfies
$\diver(f\nabla w)=-\diver(h\nabla u_0)$, $w=0$ on $\partial D$. The energy
identity (test against $w$, justified since $f\in C^{\alpha}$, $\alpha>1$
gives $w\in H^{1}$) gives
$(1-r)\norm{\nabla w}_{2}^{2}\le\int_D f\abs{\nabla w}^{2}
=-\int_D h\nabla u_0\cdot\nabla w
\le\norm{h}_2\norm{\nabla u_0}_{\infty}\norm{\nabla w}_2$, hence
$\norm{\nabla w}_2\lesssim\norm{h}_2$. Since
$R(h)=-\Delta^{-1}\diver(h\nabla w)$, for any $\psi\in L^{2}$ with
$\phi=-\Delta^{-1}\psi$,
$\inner{R(h)}{\psi}=-\int_D h\nabla w\cdot\nabla\phi$, so
$\abs{\inner{R(h)}{\psi}}\le\norm{h}_{\infty}\norm{\nabla w}_2
\norm{\nabla\phi}_2\lesssim\norm{h}_{\infty}\norm{h}_2\norm{\psi}_2$.
Duality gives the bound.
\end{proof}

\begin{remark}[$\alpha>1$ is essential]\label{rem:alpha>1}
At the minimax scale $\delta_*\asymp n^{-1/(2\alpha+d+2)}$ the wavelet
perturbations used below have $\norm{h}_{\infty}\lesssim\delta_*^{\alpha}$,
so
$\norm{R(h)}_2/(\delta_*\norm{h}_2)\lesssim\delta_*^{\alpha-1}\to0$
\emph{iff} $\alpha>1$. Thus the linearization is faithful precisely in the
regime $\alpha>1$, which is why the hypothesis appears.
\end{remark}

\section{Degree-one smoothing of the linearization}\label{sec:smoothing}

We now make Remark~\ref{rem:crux} precise with an exact Fourier-side
computation, which shows $T$ is an operator of order exactly $-1$.

\begin{lemma}[Exact smoothing; Fourier form]\label{lem:fourier}
Let $\chi\in C^{\infty}_c(\R^{d})$, $0\le\chi\le1$, $\chi=1$ on
$\SET{\abs x\le\tfrac12}$, $\supp\chi\subset\SET{\abs x\le\tfrac34}$. For
$h\in C^{\alpha}(D)$ set $\tilde h=h\chi$, extended by $0$ to $\R^{d}$.
Then, as tempered distributions,
\begin{equation}\label{eq:fourierL}
\widehat{L\tilde h}(\xi)
=-\frac{1}{d}\,\xi\cdot\nabla_{\xi}\widehat{\tilde h}(\xi).
\end{equation}
Consequently, for a \emph{scale-$\delta$ wavelet packet} $\tilde h$ and
dyadic shells $S_k=\SET{\xi:\abs\xi\in[2^{k},2^{k+1})}$,
\begin{equation}\label{eq:shell}
\norm{\,\abs\xi^{-2}\widehat{L\tilde h}(\xi)\,}_{L^{2}(S_k)}
\;=\;\frac{1}{d}\,\norm{\abs\xi^{-2}\,\xi\cdot\nabla_\xi
\widehat{\tilde h}(\xi)}_{L^{2}(S_k)},
\qquad
\norm{U\tilde h}_{L^{2}(S_k)}\asymp\delta\norm{\tilde h}_{L^{2}(S_k)}
\ \ \text{on the loaded shell }2^{k}\asymp\delta^{-1},
\end{equation}
the two-sided bound being valid \emph{for packets} (the situation in which
the lower bound is used below; it is \emph{not} a valid two-sided identity
for an arbitrary $C^\alpha$ function, whose $\widehat{\tilde h}$ can be
radially homogeneous and radially annihilate $\xi\cdot\nabla_\xi$).
\end{lemma}
\begin{proof}
Using $\widehat{x_j F}(\xi)=i\,\partial_{\xi_j}\widehat F(\xi)$ and
$\widehat{\partial_j h}(\xi)=i\xi_j\widehat h(\xi)$,
\[
\widehat{x_j\partial_j \tilde h}(\xi)
=i\,\partial_{\xi_j}\bigl(i\xi_j\widehat{\tilde h}(\xi)\bigr)
=-\partial_{\xi_j}\bigl(\xi_j\widehat{\tilde h}(\xi)\bigr)
=-\widehat{\tilde h}(\xi)-\xi_j\partial_{\xi_j}\widehat{\tilde h}(\xi).
\]
Summing over $j=1,\dots,d$:
$\widehat{A\tilde h}(\xi)=\frac1d\sum_j\widehat{x_j\partial_j\tilde h}(\xi)
=-\frac1d\bigl(d\,\widehat{\tilde h}(\xi)+\xi\cdot\nabla_\xi
\widehat{\tilde h}(\xi)\bigr)$. Hence
\[
\widehat{L\tilde h}(\xi)
=\widehat{\tilde h}(\xi)+\widehat{A\tilde h}(\xi)
=-\frac{\xi\cdot\nabla_{\xi}\widehat{\tilde h}(\xi)}{d}.
\]
This proves \eqref{eq:fourierL}. For the packet corollary \eqref{eq:shell}
of a scale-$\delta$ packet $\tilde h(x)=\sum_m\varepsilon_m a_m
\psi((x-x_m)/\delta)$, one expands
$\widehat{\tilde h}(\xi)=\delta^{d}\widehat\psi(\delta\xi)\sum_m
\varepsilon_m a_m e^{-ix_m\cdot\xi}$ and reads off the loaded-shell gain:
$\xi\cdot\nabla_\xi$ differentiates the phase $e^{-ix_m\cdot\xi}$ to bring
down a factor $x_m\cdot\xi\asymp\abs{x_m}\delta^{-1}$, and on
$\abs{x_m}\ge c_0$ this is $\asymp\delta^{-1}$, so
$\norm{\xi\cdot\nabla_\xi\widehat{\tilde h}}_{L^2(S_k)}
\asymp\delta^{-1}\norm{\widehat{\tilde h}}_{L^2(S_k)}$; dividing by
$\abs\xi^2\asymp\delta^{-2}$ gives \eqref{eq:shell} (with the phase
interference controlled by separation and Lemma~\ref{lem:ao}). The
\emph{lower} side of this packet estimate in the uniform sense relevant
for \eqref{eq:scale2} is reproved, with constants, inside
Lemma~\ref{lem:scale2}'s proof (the display labelled ``Lower bound''
there); the present form \eqref{eq:shell} records the identity content.
\end{proof}

\begin{lemma}[Almost-orthogonality of separated exponentials]\label{lem:ao}
Let $\Lambda=\SET{x_m}\subset\SET{\abs x\le\tfrac12}$ be
$4\delta$-separated, $\abs{x_m-x_{m'}}\ge4\delta$. For
$b=(b_m)\in\ell^{2}(\Lambda)$ and $\mathbf a_b(\xi):=\sum_m b_m
e^{-ix_m\cdot\xi}$, define the frequency-ball average
$\mathcal M_\rho(b):=\frac{1}{\abs{B_\rho}}\int_{B_\rho}
\abs{\mathbf a_b(\xi)}^{2}d\xi$ on $B_\rho=\SET{\abs\xi\le\rho}$,
$\rho=\Theta\delta^{-1}$, $\Theta$ a fixed constant. Then, uniformly in
$\delta$ and the configuration $\Lambda$,
\begin{equation}\label{eq:aobound}
\mathcal M_\rho(b)\ \lesssim\ \sum_m \abs{b_m}^{2},
\qquad\text{and}\qquad
\mathcal M_\rho(\SET{x_m\cdot(\cdot)}\,b)\ \lesssim\
\delta^{-2}\sum_m\abs{b_m}^{2}
\end{equation}
where the second statement applies to the weighted coefficients
$\widetilde b_m=x_m\cdot v\,b_m$ with a fixed unit vector $v$ (so
$\abs{\widetilde b_m}\le\abs{b_m}\abs{x_m}\le\tfrac12\abs{b_m}$), and the
corresponding phase function is $(x_m\cdot v)e^{-ix_m\cdot\xi}$.
\end{lemma}
\begin{proof}
Write, with $\chi_\rho=\one_{B_\rho}$ and
$\checkK{\chi_\rho}$ its inverse Fourier transform,
$\int_{B_\rho}\abs{\mathbf a_b}^{2}d\xi
=\sum_{m,m'}b_m\overline{b_{m'}}\,\checkK{\chi_\rho}(x_{m'}-x_m)$. The
diagonal terms ($m=m'$) give
$\sum_m\abs{b_m}^2\checkK{\chi_\rho}(0)=\abs{B_\rho}\sum_m\abs{b_m}^2$
since $\checkK{\chi_\rho}(0)=\abs{B_\rho}$. For the off-diagonal part we
cannot use the ball indicator's transform (a jinc/Bessel profile decaying
only as $|y|^{-(d+1)/2}$), whose Schur sum diverges; instead dominate by a
\emph{smooth majorant}: choose $\varphi_0\in C_c^\infty$ with
$\varphi_0(\eta)\ge\one_{|\eta|\le1}$ and
$\abs{\checkK{\varphi_0}(y)}\lesssim(1+\abs y)^{-N}$ for any $N$ (a
Schwartz function since $\varphi_0$ is smooth compactly supported), set
$\varphi_\rho(\xi)=\varphi_0(\xi/\rho)$, and note
$\checkK{\varphi_\rho}(y)=\rho^{d}\checkK{\varphi_0}(\rho y)$ with
$\checkK{\varphi_\rho}(0)=\rho^{d}\checkK{\varphi_0}(0)\asymp\abs{B_\rho}$
and $\abs{\checkK{\varphi_\rho}(y)}\lesssim\rho^{d}(1+\rho\abs{y})^{-N}$.
Because $\one_{B_\rho}\le\varphi_\rho$,
$\int_{B_\rho}\abs{\mathbf a_b}^{2}d\xi
\le\int\varphi_\rho\abs{\mathbf a_b}^{2}d\xi
=\sum_{m,m'}b_m\overline{b_{m'}}\checkK{\varphi_\rho}(x_{m'}-x_m)$, and the
off-diagonal part is bounded by
\[
\sum_{m\ne m'}\abs{b_m}\abs{b_{m'}}\,\rho^{d}
\bigl(1+\rho\abs{x_{m'}-x_m}\bigr)^{-N}
\lesssim\rho^{d}\sum_{m,m'}\frac{\abs{b_m}\abs{b_{m'}}}
{(1+\abs{x_{m'}-x_m}/\delta)^{N}},
\]
using $\rho= \Theta\delta^{-1}$ and $\abs{x_{m'}-x_m}\ge4\delta$. The kernel
$K(x_{m'},x_m)=(1+\abs{x_{m'}-x_m}/\delta)^{-N}$ on the separated set is
Schur-/Cotlar-summable: $\sup_{m'}\sum_m K\lesssim\sum_{k\ge1}k^{d-1}
(1+4k)^{-N}<\infty$ for $N>d$, and symmetrically in $m'$; hence the
off-diagonal quadratic form is $\lesssim\rho^{d}\sum_m\abs{b_m}^{2}$.
Dividing by $\abs{B_\rho}\asymp\rho^{d}$ gives the first bound of
\eqref{eq:aobound}. The second (weighted) bound is the first applied to the
coefficients $b_m^{(j)}=a_m x_{m,j}$ for each coordinate $j$: pointwise,
$\abs{\sum_m a_m(x_m\cdot\xi)e^{ix_m\cdot\xi}}^{2}
=\abs{\sum_j\xi_j\sum_m a_m x_{m,j}e^{ix_m\cdot\xi}}^{2}
\le\abs{\xi}^{2}\sum_j\abs{\sum_m a_m x_{m,j}e^{ix_m\cdot\xi}}^{2}$, and
averaging each coordinate term over $B_\rho$ with (i), using
$\abs{x_{m,j}}\le\abs{x_m}\le\tfrac12$ and $\abs\xi\le\rho$,
\[
\frac{1}{\abs{B_\rho}}\int_{B_\rho}
\Bigl|\sum_m a_m(x_m\cdot\xi)e^{ix_m\cdot\xi}\Bigr|^{2}d\xi
\le\rho^{2}\sum_j\mathcal M_\rho(a\,x_{\cdot,j})
\lesssim\rho^{2}\sum_j\sum_m a_m^{2}x_{m,j}^{2}
\le\frac{\rho^{2}}{4}\sum_m a_m^{2}
\asymp\delta^{-2}\sum_m a_m^{2}.
\]
This proves the second bound of \eqref{eq:aobound}.
\end{proof}

\begin{lemma}[Scale estimate for wavelet packets]\label{lem:scale2}
Fix $\chi$ as above and $R\in(0,\tfrac18]$. There are $c_1,C_1>0$
(depending only on $d,\chi,R$) such that the following holds. Let
$g=\sum_{m}\varepsilon_m a_m\psi\bigl(\tfrac{\cdot-x_m}{\delta}\bigr)$ be a
wavelet packet built from a fixed $\psi\in C^\infty_c$ with
$\supp\psi\subset B(0,R)$, $\psi\ne0$, $\int\psi\,dx=0$, where
$\varepsilon_m\in\{-1,0,1\}$, $a_m\in\R$, and the centres $x_m$ lie in the
\emph{annulus} $c_0\le\abs{x_m}\le\tfrac12-c_0$ for a fixed constant
$c_0=c_0(d)\in(0,\tfrac{1}{16}]$ with pairwise separation
$\abs{x_m-x_{m'}}\ge4\delta$ (so the terms have pairwise disjoint supports,
all inside $\SET{\abs x\le\tfrac12-c_0/2}$, at distance
$\ge c_0/2+R\,\delta\ge c_0/2$ from $\partial D$; the support stays in
$\SET{\abs x\le\tfrac34}$ too). Then for $\delta\in(0,\delta_0]$,
\begin{equation}\label{eq:scale2}
c_1\,\delta\,\norm{g}_{L^{2}(D)}
\ \le\ \norm{Tg}_{L^{2}(D)}
\ \le\ C_1\,\delta\,\norm{g}_{L^{2}(D)},
\end{equation}
with $\norm{g}_2$ computed on $D$ (equivalently $\R^d$, by the support).
\end{lemma}
\begin{proof}
Since $\supp g\subset\SET{\abs x\le\tfrac34}$ and $\chi=1$ there, $g=g\chi$,
so $\tilde g=g$. Write
$Tg=U-q$ with $U=-\Delta^{-1}_{\R^d}(Lg)$ and $q$ harmonic in $D$,
$q=U$ on $\partial D$.

\smallskip
\noindent\textbf{Whole-space term $U$.} We use the exact identity
\eqref{eq:fourierL}: with $\Lambda F=\xi\cdot\nabla_\xi F$ the dilation
generator, $U=\frac1d\abs D^{-2}\Lambda L_{\rm op}\,g$ is a nonlocal
Fourier operator of order exactly $-1$, and on a scale-$\delta$ packet it
multiplies the $L^{2}$ norm by $\asymp\delta$ (the content of
\eqref{eq:shell}). For concreteness and constants we compute directly.
Write
$\widehat g(\xi)=\delta^{d}\widehat\psi(\delta\xi)\mathbf a(\xi)$ with
$\mathbf a(\xi):=\sum_m\varepsilon_m a_m e^{-ix_m\cdot\xi}$. Since the
packet terms are disjointly supported,
$\norm{g}_2^{2}=\sum_m a_m^{2}\delta^{d}\norm{\psi}_2^{2}$. With
dilation generator $\Lambda F=\xi\cdot\nabla_\xi F$ applied to
$F=\widehat g$, and splitting,
$\xi\cdot\nabla_\xi\widehat g
=\delta^{d}\bigl(B_1(\xi)+B_2(\xi)\bigr)$ with
$B_1(\xi)=(\Lambda\widehat\psi)(\delta\xi)\,\mathbf a(\xi)$,
$B_2(\xi)=\widehat\psi(\delta\xi)\,\xi\cdot\nabla_\xi\mathbf a(\xi)$, and
$\xi\cdot\nabla_\xi\mathbf a(\xi)
=-i\sum_m\varepsilon_m a_m (x_m\cdot\xi)e^{-ix_m\cdot\xi}$.
Cover the loaded shell $S_k=\SET{\abs\xi\asymp\delta^{-1}}$ by $O(1)$
balls $B_\rho$, $\rho\asymp\delta^{-1}$, and apply Lemma~\ref{lem:ao}
(claims (i) and the weighted corollary):
$\frac1{\abs{B_\rho}}\int_{B_\rho}\abs{\mathbf a}^{2}\lesssim\sum a_m^{2}$
and
$\frac1{\abs{B_\rho}}\int_{B_\rho}\abs{\xi\cdot\nabla_\xi\mathbf a}^{2}
\lesssim\rho^{2}\cdot\tfrac14\sum a_m^{2}\asymp\delta^{-2}\sum a_m^{2}$.
Since $\abs{\widehat\psi(\delta\xi)},\abs{(\Lambda\widehat\psi)(\delta\xi)}
\lesssim$ const, squaring and integrating gives, per contributing term,
\begin{align}\label{eq:uterm}
\int_{S_k}\abs\xi^{-4}\abs{\xi\cdot\nabla_\xi\widehat g}^{2}d\xi
&\lesssim\delta^{2d+4}\Bigl[
\int_{S_k}\abs{B_1}^{2}+\int_{S_k}\abs{B_2}^{2}\Bigr]
\nonumber\\
&\lesssim\delta^{2d+4}\cdot\abs{S_k}\,
\frac1{\abs{B_\rho}}\Bigl[
\int_{B_\rho}\abs{\mathbf a}^{2}
+\int_{B_\rho}\abs{\xi\cdot\nabla_\xi\mathbf a}^{2}\Bigr]
\nonumber\\
&\lesssim\delta^{2d+4}\cdot\delta^{-d}
\Bigl(\sum a_m^{2}+\delta^{-2}\sum a_m^{2}\Bigr)
\ \lesssim\ \delta^{2}\cdot\delta^{d}\sum a_m^{2}
\ =\ C\,\delta^{2}\,\frac{\norm{g}_2^{2}}{\norm{\psi}_2^{2}} .
\end{align}
Here the power count is: $\delta^{2d}$ (prefactor) $\times\delta^{4}$
($|\xi|^{-4}$) $\times\delta^{-d}$ (shell volume) $\times\delta^{-2}$
(weighted AO, dominant $B_2$ term) $=\delta^{2}$; the $B_1$ term is one
power of $\delta$ smaller. The constant is uniform in $M$ since there are
$O(1)$ balls.

For the off-shell shells $S_j$: those with $2^j\gg\delta^{-1}$ use the
Schwartz decay $\abs{\widehat\psi(\delta\xi)}\lesssim(\delta\abs\xi)^{-N}$
so each contributes $\delta^{2}\cdot 2^{-N(j-k)}$; those with
$2^j\ll\delta^{-1}$ use the mean-zero bound
$\abs{\widehat\psi(\delta\xi)}\lesssim\delta\abs\xi$
($\int\psi=0$), contributing $\delta^{2}\cdot 2^{-2(k-j)}$. Both tails sum
geometrically, so the whole-space term satisfies the upper bound
\begin{equation}\label{eq:Uupper}
\norm{U}_{L^{2}(\R^{d})}^{2}\ \lesssim\ \delta^{2}\norm{g}_{2}^{2},
\qquad\text{uniformly in }M,\delta .
\end{equation}

\noindent\textbf{Lower bound (loaded-shell diagonal dominance).} Retain
only the loaded shell $S_k=\SET{\abs\xi\asymp\delta^{-1}}$ (all terms
non-negative on the Fourier side). Since
$\norm{U}_2^{2}\ge\int_{S_k}\abs\xi^{-4}\abs{\widehat{Lg}}^{2}d\xi
\asymp\delta^{4}\int_{S_k}\abs{\widehat{Lg}}^2d\xi$ and by
\eqref{eq:fourierL} $\widehat{Lg}=-\frac1d\xi\cdot\nabla_\xi\widehat g$,
using $\abs{a+b}^{2}\ge\tfrac12\abs{b}^{2}-\abs{a}^{2}$ with
$b=\delta^{d}\widehat\psi(\delta\xi)\xi\cdot\nabla_\xi\mathbf a(\xi)$ and
$a=\delta^{d}(\Lambda\widehat\psi)(\delta\xi)\mathbf a(\xi)$:
\[
\norm{U}_2^{2}
\ge\frac{\delta^{4}}{2d^{2}}\int_{S_k}\delta^{2d}
\abs{\widehat\psi(\delta\xi)}^{2}\abs{\xi\cdot\nabla_\xi\mathbf a(\xi)}^{2}d\xi
-\frac{\delta^{4}}{d^{2}}\int_{S_k}\delta^{2d}
\abs{(\Lambda\widehat\psi)(\delta\xi)}^{2}\abs{\mathbf a(\xi)}^{2}d\xi .
\]
Choose once and for all the base profile $\psi$ of the lemma
statement: a radial $\psi\in C^\infty_c$ with $\supp\psi\subset B(0,R)$,
$\int\psi=\widehat\psi(0)=0$, and $\widehat\psi$ nonvanishing on the
closed annulus $\SET{1\le\abs\eta\le2}$ (readily constructed by applying a
power of $-\Delta$ to a radial $C^\infty_c$-bump and adding a small
correction). Thus $\abs{\widehat\psi(\delta\xi)}\ge c_\psi>0$ on the
loaded shell while $\abs{(\Lambda\widehat\psi)(\delta\xi)}\le C_\psi$. Now
$\xi\cdot\nabla_\xi\mathbf a(\xi)=-i\sum_m\varepsilon_m a_m(x_m\cdot\xi)
e^{-ix_m\cdot\delta\xi}$; because the centres lie in the annulus
$\abs{x_m}\ge c_0$, on the shell
$\abs{x_m\cdot\xi}\asymp\abs{\xi}\abs{x_m}\ge\tfrac{c_0}{2}\delta^{-1}$ (in
the diagonal/phase-gradient part). Apply Lemma~\ref{lem:ao}'s weighted
corollary with the \emph{lower} direction: the oscillatory off-diagonals
are Schur-summable (by the smooth majorant bound of the lemma) and are
dominated by the diagonal, so
\[
\int_{S_k}\abs{\widehat\psi(\delta\xi)}^{2}
\abs{\xi\cdot\nabla_\xi\mathbf a(\xi)}^{2}d\xi
\ \ge\ c_\psi^{2}\,\delta^{-d}\,\delta^{-2}\,\frac{c_0^{2}}{4}\sum_m a_m^{2}
\ \ \bigl(\text{diagonal }m=m'\bigr),
\]
the $m=m'$ terms giving
$\sum_m a_m^{2}\int_{S_k}\abs{x_m\cdot\xi}^{2}d\xi
\asymp\sum_m a_m^{2}\delta^{-d}\delta^{-2}$ since
$\int_{S_k}\abs{x_m\cdot\xi}^{2}d\xi\asymp\abs{x_m}^{2}\delta^{-2-d}\cdot
\delta^{2}=\abs{x_m}^{2}\delta^{-d}\gtrsim c_0^{2}\delta^{-d}$.
Similarly the second (subtracted) term is bounded above by the
\emph{unweighted} AO: $\le C_\psi\delta^{-d}\sum a_m^{2}$ (no
$\delta^{-2}$). Hence
\[
\norm{Ug}_2^{2}
\ \ge\ \delta^{4+2d}\,\delta^{-d}\Bigl(
\frac{c_\psi^{2}c_0^{2}}{8d^{2}}\,\delta^{-2}\sum a_m^{2}
-\frac{C_\psi}{d^{2}}\sum a_m^{2}\Bigr)
\ \ge\ c\,\delta^{d+4}\cdot\delta^{-2}\sum a_m^{2}
\ =\ c\,\delta^{2}\cdot\delta^{d}\sum a_m^{2}
\ =\ c\,\delta^{2}\,\frac{\norm{g}_2^{2}}{\norm{\psi}_2^{2}},
\]
using $\delta^{4+2d}\delta^{-d}\delta^{-2}=\delta^{d+2}
=\delta^{2}\cdot\delta^{d}$, and for $\delta\le\delta_0(c_\psi,c_0,C_\psi)$
the dominant $\delta^{-2}$ term controls the subtracted one. Hence
$\norm{U}_2\ge c_U'\,\delta\norm{g}_2$ with
$c_U'=c_U'(d,\psi,c_0)>0$.

\smallskip
\noindent\textbf{Harmonic correction $q$.} By Lemma~\ref{lem:harmonic}
(applied with $F=Lg=\frac1d\diver(xg)$, for which the solution $w=Tg$
splits as $w=U-q$ with exactly the present $U$ and $q$),
$U=(-\Delta_{\R^d})^{-1}F$ is a convolution with a locally-integrable
kernel of a field supported in $D_{(\tau_0)}$, $\tau_0=\tfrac18$, and on
$\partial D$ one has the interior bound
$\norm{U}_{L^{\infty}(\partial D)}\lesssim\norm g_2$
(Cauchy--Schwarz in $y$, distance $\ge\tau_0/4$ to the support). Split
$q=q_0+q_{\rm osc}$ with $q_0=\abs{\partial D}^{-1}\int_{\partial D}q$ the constant Poisson
mode; the maximum principle gives $\norm{q_0}_{L^{2}}
\le\abs{\partial D}^{-1}\int_{\partial D}\abs U\lesssim\delta\norm g_2$ by the bump-wise
multipole expansion of $U|_{\partial D}$ (monopole killed by
$\int\psi=0$, first-order remainder $O(\delta\norm{\psi_m}_1)$,
$\sum_m\abs{a_m}\norm{\psi_m}_1\lesssim\norm g_2$), and the
mean-zero part decays into the interior at the Poisson rate:
$q_{\rm osc}=\sum_{\Lambda}\rho^{\sqrt\Lambda}\widehat U(\Lambda)$, so
on $\SET{\abs x\le\tfrac34}$ (which carries the shell mass of $U$)
$\norm{q_{\rm osc}}_{L^{2}}\le\bigl(\max_\Lambda
\rho^{2\sqrt\Lambda}\bigr)^{1/2}\norm U_{L^{2}(\partial D)}
\lesssim e^{-c/\delta}\norm g_2\ll\delta\norm g_2$ (the largest
$\sqrt\Lambda$ is $\asymp\delta^{-1}$, and
$\rho^{2\sqrt\Lambda}\le(\frac34)^{2c\delta^{-1}}\le e^{-c'/\delta}$).
Hence
\begin{equation}\label{eq:qbound}
\norm{q}_{L^{2}(\abs x\le\frac34)}\ \lesssim\ \delta\,\norm g_2 ,
\qquad\text{constant depending on }(d,\psi,c_0,\tau_0),
\end{equation}
$\SET{\abs x\le\tfrac34}$ being where the shell mass of $U$ lies
($\supp g\subset\SET{\abs x\le\tfrac34}$, the $\R^d$ integral
dominated by its $\SET{\abs x\le\tfrac34}$ part up to the
$e^{-c/\delta}$ tail already accounted).

\smallskip
\noindent Combining,
$\norm{Tg}_2\le\norm{U}_2+\norm{q}_2\lesssim\delta\norm g_2$ and
$\norm{Tg}_2\ge\norm{U}_{L^{2}(\abs x\le\frac34)}-\norm{q}_{L^{2}(\abs
x\le\frac34)}\ge(c_U'/2-C_q\delta)\,\delta\norm g_2$ for the
$\SET{\abs x\le\frac34}$-restricted shell lower bound $c_U'/2$ (same
argument, all estimates on the shrunk domain) and $C_q$ the constant of
\eqref{eq:qbound}; with $\delta_0$ shrunk so $C_q\delta_0\le c_U'/4$,
\[
\norm{Tg}_2\ \ge\ \tfrac{c_U'}4\,\delta\norm g_2 =:c_1\delta\norm g_2,
\qquad
\norm{Tg}_2\ \le\ (C_U+C_q)\delta\norm g_2=:C_1\delta\norm g_2 ,
\]
with $c_1,C_1,\delta_0$ depending only on $(d,\psi,c_0)$. (The
origin-centre modes, excluded here by the annulus hypothesis, are
treated in Lemma~\ref{lem:invert} via the exact radial computation.)
\end{proof}

\begin{remark}[Constant bookkeeping in \eqref{eq:scale2}]\label{rem:const}
Both $C_E''=C_E''(d)$ (trace/elliptic constant of the unit ball) and
$c_U'=c_U'(d,\psi,c_0)$ are $(\alpha,n,\delta,r)$-independent, so the
threshold $\delta_0=\delta_0(c_\psi,c_0,C_\psi,C_E'')$ and all constants
below keep the $(\alpha,d,\psi)$-dependence only; boundary-near bumps are
\emph{not} asked to be small, the loss being absorbed by shrinking
$\delta_0$.
\end{remark}

\begin{remark}[$T$ is a degree-one operator, globally]\label{rem:degree1}
For a full wavelet expansion $h=\sum_{j\ge j_0}h_j$ with $h_j$ at scale
$2^{-j}$, applying Lemma~\ref{lem:scale2} to each dyadic piece and summing
(the pieces are almost orthogonal in $L^2$, and $T$ maps the scale
$2^{-j}$ piece into functions still almost-orthogonal up to the operator's
pseudo-local tail) yields
$\norm{Th}_2^{2}\lesssim\sum_j 2^{-2j}\norm{h_j}_2^{2}
\lesssim\norm{h}_{H^{-1}}^{2}$, i.e.\ $T$ maps $H^{s}$ into $H^{s+1}$ type
spaces with exactly one derivative of smoothing, globally on the Besov
scale. In particular $\norm{Th}_2\le C\norm{h}_2$, the only global fact
(besides \eqref{eq:scale2}) used below.
\end{remark}

\section{Lower bound}\label{sec:lower}

\begin{proposition}\label{prop:lower}
There exist $r_0=r_0(\alpha,d)\in(0,1]$ and $c>0$ such that for all
$r\in(0,r_0]$ and all large $n$,
\[
\inf_{\widehat f_n}\sup_{f_0\in\F(\alpha,r)}
\E_{f_0}\norm{\widehat f_n-f_0}_{2}^{2}\ \ge\ c\,n^{-\frac{2\alpha}{2\alpha+d+2}} .
\]
\end{proposition}

\begin{proof}
\textbf{Step 1 (wavelet packet).} Fix the single base profile
$\psi\in C^\infty_c(\R^{d})$ of Lemma~\ref{lem:scale2} (radial,
$\supp\psi\subset\SET{\abs x\le\tfrac14}$, $\int\psi=\widehat\psi(0)=0$,
$\widehat\psi$ nonvanishing on $\SET{1\le\abs\eta\le2}$), normalized so
$\norm{\psi}_{C^\alpha}=1$; then $\norm{\psi}_2>0$ and
$\norm\psi_{C^\alpha}=1$ are fixed $(\alpha,d)$-constants. For $\delta\in(0,\delta_0]$ choose a maximal
collection $\SET{x_m}_{m=1}^{M}\subset\SET{c_0\le\abs{x}\le\tfrac12}$ with
pairwise distance $\ge4\delta$ (hence $M\asymp\delta^{-d}$), and set
$\psi_m(x)=\psi\bigl(\tfrac{x-x_m}{\delta}\bigr)$. Then
$\supp\psi_m\subset B(x_m,\delta/4)$,
$\norm{\psi_m}_\infty=\norm{\psi}_\infty$,
$\norm{\psi_m}_2=\delta^{d/2}\norm{\psi}_2$,
$\norm{\psi_m}_{C^\alpha}=\delta^{-\alpha}\norm{\psi}_{C^\alpha}$ (the
$C^\alpha$ norm is $\asymp\delta^{-\alpha}$ for $\delta\le1$: derivatives
scale as $\delta^{-j}$ and the top $t$-H\"older/Zygmund seminorm as
$\delta^{-\alpha}$, dominating the lower-order terms).
The supports of the $\psi_m$ are pairwise disjoint and contained in
$\SET{\abs x\le\tfrac34}$, so every partial sum
$\sum_{m\in S}\pm\psi_m$ is a wavelet packet of the form treated in
Lemma~\ref{lem:scale2} (fixed base profile $\psi$, centres in
$\SET{c_0\le\abs{x}\le\tfrac12}$, separation $\ge4\delta$; the annulus
restriction keeps the two-sided scale estimate of
Lemma~\ref{lem:scale2}).

\textbf{Step 2 (hypotheses).} By Varshamov--Gilbert, there is a family
$\mathcal S$ of subsets $S\subset\SET{1,\dots,M}$ with $\abs S=M/2$, each
pair $S\ne S'$ at Hamming distance $\abs{S\triangle S'}\ge M/8$, and
$\log\abs{\mathcal S}\ge c_0 M$ for an absolute $c_0>0$. Index hypotheses
by $S\in\mathcal S$ and a sign/scale parameter $\kappa>0$ (fixed below):
\[
f_S=\one+\kappa\delta^{\alpha}\Psi_S,
\qquad\Psi_S:=\sum_{m\in S}\psi_m .
\]
Because the $\psi_m$ have disjoint supports with the same profile,
$\norm{\Psi_S}_{C^\alpha}\lesssim\max_{m}\norm{\psi_m}_{C^\alpha}
\asymp\delta^{-\alpha}\norm{\psi}_{C^\alpha}$: the $C^a$ norm of the sum is
the max over the summands (disjoint supports), and the top
$t$-H\"older/Zygmund difference quotient satisfies
$\sup_{x,y}\abs{\nabla^{a}\Psi_S(x)-\nabla^{a}\Psi_S(y)}/\abs{x-y}^{t}
\le 2^{\,1-t}\delta^{-\alpha}\norm{\nabla^{a}\psi}_{\infty}
+\delta^{-\alpha}[\psi]_{\dot C^{t}}$ — cross-support quotients are the
same order as single-bump ones, giving comparability with a
$\psi$-dependent constant $\le 2^{1-t}+1$. Hence
\[
\norm{f_S-\one}_{C^\alpha}
\lesssim\kappa\delta^{\alpha}\cdot\delta^{-\alpha}\norm{\psi}_{C^\alpha}
=\kappa\norm{\psi}_{C^\alpha}\ \le\ r
\]
provided $\kappa\le r\norm{\psi}_{C^\alpha}^{-1}$, i.e.\ $\kappa\le r$
since $\norm{\psi}_{C^\alpha}=1$ (Step~1). We take the maximal admissible
$\kappa=r$; the separation below is then $\varepsilon^{2}\asymp
r^{2}\delta^{2\alpha}$, so the lower-bound constant inherits the factor
$r^{2}$ of the shrinking class $\F(\alpha,r)$ -- a fixed positive constant
for every fixed $r\in(0,r_0]$. Thus $f_S\in\F(\alpha,r)$ for all
$S\in\mathcal S$. Also, by disjoint supports,
\[
\norm{\Psi_S}_2^{2}=\abs S\,\norm{\psi_m}_2^{2}\asymp M\,\delta^{d}\asymp1,
\qquad
\norm{\Psi_S}_{\infty}=\norm{\psi}_{\infty}.
\]

\textbf{Step 3 (separation).} For $S\ne S'$,
$f_S-f_{S'}=\kappa\delta^{\alpha}\sum_{m\in S\triangle S'}\pm\psi_m$, the
supports disjoint, so
\[
\norm{f_S-f_{S'}}_2^{2}
=\kappa^{2}\delta^{2\alpha}\abs{S\triangle S'}\,\norm{\psi_m}_2^{2}
\asymp\kappa^{2}\delta^{2\alpha}\cdot M\cdot\delta^{d}
\asymp\kappa^{2}\delta^{2\alpha}.
\]
So the family is $\varepsilon$-separated with
$\varepsilon^{2}\asymp\kappa^{2}\delta^{2\alpha}$.

\textbf{Step 4 (KL control).} With $g_S:=f_S-\one=\kappa\delta^{\alpha}\Psi_S$:
$g_S$ is supported in $\SET{\abs x\le\tfrac34}$, is a single-scale packet
of the form in Lemma~\ref{lem:scale2} (fixed profile $\psi$ with
$\int\psi=0$, centres in $\SET{\abs x\le\tfrac12}$ separated by
$\ge4\delta$; each $\psi_m$ has
$\norm{\nabla\psi_m}_2\asymp\delta^{-1}\norm{\psi_m}_2$), so by
Lemma~\ref{lem:scale2},
\[
\norm{Tg_S}_2\ \asymp\ \delta\norm{g_S}_2
\ \asymp\ \delta\cdot\kappa\delta^{\alpha}\norm{\Psi_S}_2
\ \asymp\ \kappa\delta^{\alpha+1}.
\]
(Here $\norm{\Psi_S}_2\asymp1$ by Step 2.) The quadratic remainder, by
Lemma~\ref{lem:quad},
\[
\norm{R(g_S)}_2\lesssim\norm{g_S}_\infty\norm{g_S}_2
=\kappa\delta^{\alpha}\norm{\psi}_\infty\cdot\kappa\delta^{\alpha}\norm{\Psi_S}_2
\asymp\kappa^{2}\delta^{2\alpha}.
\]
Since $\alpha>1$, $\norm{R(g_S)}_2\lesssim\kappa^{2}\delta^{2\alpha}
\ll\kappa\delta^{\alpha+1}\asymp\norm{Tg_S}_2$ for $\delta$ small (the
ratio is $\kappa\delta^{\alpha-1}\to0$ --- the decisive use of $\alpha>1$).
Therefore
\[
\norm{u_{f_S}-u_0}_2
\le\norm{Tg_S}_2+\norm{R(g_S)}_2
\lesssim\kappa\delta^{\alpha+1},
\]
and by \eqref{eq:KL},
\begin{equation}\label{eq:KLbound}
\operatorname{KL}(P_{f_S}\,\|\,P_{f_0})
=\frac{n}{2}\norm{u_{f_S}-u_0}_2^{2}
\le Cn\kappa^{2}\delta^{2\alpha+2}.
\end{equation}

\textbf{Step 5 (Fano).} Take
$\delta=c\,n^{-1/(2\alpha+d+2)}$. Then the balance is
\[
n\delta^{2\alpha+2}=c^{2\alpha+2}n^{1-\frac{2\alpha+2}{2\alpha+d+2}}
=c^{2\alpha+2}n^{\frac{d}{2\alpha+d+2}}
=c^{2\alpha+2}\delta^{-d}.
\]
Choosing $c>0$ small makes $Cn\kappa^{2}\delta^{2\alpha+2}
\le\frac1{10}\,\delta^{-d}\le\frac1{10}\cdot\frac{\log\abs{\mathcal S}}{c_0}$,
i.e.\ the average KL is $\le\frac1{10}\log\abs{\mathcal S}$. With
$\varepsilon^{2}\asymp\kappa^{2}\delta^{2\alpha}
=\kappa^{2}c^{2\alpha}n^{-\frac{2\alpha}{2\alpha+d+2}}$ and
$\log\abs{\mathcal S}\ge c_0M\asymp\delta^{-d}$, Fano's inequality
(Lemma~\ref{lem:fano}) gives
\[
\max_{S\in\mathcal S}\E_{f_S}\norm{\widehat f-f_S}_2^{2}
\ \ge\ \frac{\varepsilon^{2}}{8}\Bigl(1-\frac1{10}-\frac{\log2}{\log\abs{\mathcal S}}\Bigr)
\ \ge\ c'\,n^{-\frac{2\alpha}{2\alpha+d+2}}
\]
for $n$ large. Since $\SET{f_S}\subset\F(\alpha,r)$, the minimax risk over
$\F(\alpha,r)$ is at least this, proving Proposition~\ref{prop:lower}.
\end{proof}

\begin{lemma}[Fano, applied form]\label{lem:fano}
Let $\abs{\Theta}=N\ge2$, $\SET{P_\theta}_{\theta\in\Theta}$ probability
measures with $\frac1N\sum_\theta\operatorname{KL}(P_\theta\|P_0)\le\beta\log N$,
$\beta<1$, and parameter values separated:
$\norm{f_\theta-f_{\theta'}}_2\ge\varepsilon$ for $\theta\ne\theta'$. Then
\[
\inf_{\widehat f}\max_{\theta}\E_{\theta}\norm{\widehat f-f_\theta}_2^{2}
\ \ge\ \frac{\varepsilon^{2}}{4}\Bigl(1-\beta-\frac{\log2}{\log N}\Bigr)\cdot\frac1{2}
\ \ge\ \frac{\varepsilon^{2}}{8}\Bigl(1-\beta-\frac{\log2}{\log N}\Bigr).
\]
\end{lemma}
\begin{proof}
Given any estimator $\widehat f$, let
$\widehat\theta=\arg\min_{\theta\in\Theta}\norm{\widehat f-f_\theta}_2$. By
$\varepsilon$-separation and the triangle inequality,
$\norm{\widehat f-f_\theta}_2\ge\varepsilon/2$ whenever
$\widehat\theta\ne\theta$. Markov gives
$\E_\theta\norm{\widehat f-f_\theta}_2^{2}\ge\frac{\varepsilon^{2}}{4}
P_\theta(\widehat\theta\ne\theta)$. Averaging and using
$\max_\theta\ge\frac1N\sum_\theta$:
$\max_\theta\E_\theta\norm{\widehat f-f_\theta}_2^{2}
\ge\frac{\varepsilon^{2}}{4}\cdot\frac1N\sum_\theta
P_\theta(\widehat\theta\ne\theta)$. Fano's inequality bounds the average
testing error below by
$1-\frac{\beta\log N+\log2}{\log N}=\bigl(1-\beta-\frac{\log2}{\log N}\bigr)$.
Combining gives the claim (Tsybakov, \emph{Introduction to Nonparametric
Estimation}, Theorem 2.5).
\end{proof}

\section{Upper bound}\label{sec:upper}

\begin{proposition}\label{prop:upper}
There is an estimator $\widehat f_n$ with
\[
\sup_{f_0\in\F(\alpha,r)}
\E_{f_0}\norm{\widehat f_n-f_0}_{2}^{2}
\ \le\ C\,n^{-\frac{2\alpha}{2\alpha+d+2}} .
\]
\end{proposition}

The construction: (A) discretize the observation and set up a
scale-$\delta$ sieve least-squares problem; (B) a discrete
inf--sup/smoothing inequality inverting $T$ on the sieve; (C) balance
bias, variance, and the quadratic remainder.

\subsection{(A) Observation discretization; sieve least squares}
Let $V_\delta\subset L^{2}(D)$ be the span of the Dirichlet eigenfunctions
of $-\Delta$ with eigenvalue $\le C_W\,\delta^{-2}$ ($C_W$ a fixed
large constant chosen in Lemma~\ref{lem:invert}), i.e.\ the eigenfunctions
oscillating at frequencies $\lesssim\delta^{-1}$; by Weyl's law
$\dim V_\delta\asymp\delta^{-d}$ (wavelength-$2^{-j}$ eigenfunctions carry
eigenvalues $\asymp2^{2j}$, so the cutoff $\lesssim\delta^{-2}$ keeps
\emph{physical scale $\gtrsim\delta$}). Let $P_\delta$ be the orthogonal
projection onto $V_\delta$ and set
\[
Y_\delta:=P_\delta Y=P_\delta u_{f_0}+n^{-1/2}\xi_\delta,
\qquad
\xi_\delta:=P_\delta\W,\quad \E\norm{\xi_\delta}_2^{2}=\dim V_\delta\asymp\delta^{-d}.
\]
The sieve must avoid both degeneracies of the transport field
$\nabla u_0=x/d$: the origin (where it vanishes) and the boundary
collar (where the H\"older boundary estimate loses one degree). We
therefore estimate only interior content. Let $\F_\delta$ be the span
of all wavelet packets at scales $\ge\delta$ \emph{whose centres lie
in the fixed interior} $D_1:=\SET{c_1\le\abs{x}\le\tfrac12}$ for a small
fixed $c_1=c_1(d)>0$, intersected with $\F(\alpha,r)$. Fix a patch
$\zeta$, $\equiv1$ on $\operatorname{supp}(\F_\delta-\one)$,
$\equiv0$ on $B_{c_1/2}\cup\SET{\abs x>\tfrac34}$. For a general
$f\in\F(\alpha,r)$ write $f=f_{\le\delta}+f_{>\delta}$ (scales
$\ge\delta$ and $<\delta$; $\supp(f_{\le\delta}-\one)
\subset\SET{\zeta=1}$) and set
\begin{equation}\label{eq:Pidelta}
\Pi_\delta f\ :=\ \one+\bigl(f_{\le\delta}-\one\bigr)\,\zeta ,
\qquad\text{so }\Pi_\delta f=\one \text{ off }\operatorname{supp}\zeta\ (\text{a ball about }0\text{ and a collar of }\partial D).
\end{equation}
Thus $\Pi_\delta f\in\F(\alpha,2r)$, and the comparison point
$m=\Pi_\delta f_0$ is Dirichlet ($u_m=0$ on $\partial D$, as is
$u_{f_0}$, by \eqref{eq:veryweak}) and interior-regular. By the standard
$B^\alpha_{\infty,\infty}=C^\alpha$ wavelet characterization,
\begin{equation}\label{eq:approx}
\norm{f_{>\delta}}_{\infty}\le C\delta^{\alpha},
\quad
\norm{f_{>\delta}}_{2}\le C\delta^{\alpha},
\quad
\norm{f_{\le\delta}}_{C^\alpha}\le C\norm{f}_{C^\alpha},
\quad
\dim\F_\delta\lesssim\delta^{-d}.
\end{equation}
The \emph{bias} $e=\Pi_\delta f-f=-f_{>\delta}+(f_{\le\delta}-\one)
(\zeta-\one)$ has $\norm{e}_2\lesssim\delta^{\alpha}$ by
\eqref{eq:approx}; its dominant content is interior (the
origin/boundary remainder is handled in the collar split of
\eqref{eq:finebias}). The projection $P_\delta$ retains the part of
$u_{f_0}$ oscillating at scales $\gtrsim\delta$ and makes the noise
finite-dimensional, so the estimator and its risk are well-defined on
the white-noise model. Define the sieve
least-squares estimator
\[
\widehat f=\arg\min_{g\in\F_\delta}
\ \norm{P_\delta u_g-Y_\delta}_{2}^{2}.
\]
Since $\F_\delta$ is closed and
$g\mapsto\norm{P_\delta u_g-Y_\delta}_2$ is continuous and coercive
($\norm{P_\delta u_g}_2\to\infty$ with $\norm g_2$ on the sieve by
\eqref{eq:infsup} and the apriori bound, Lemma~\ref{lem:apriori}), a
measurable minimizer $\widehat f$ exists (selection theorem on the
closed finite-dimensional sieve); only its risk is analyzed below.

\subsection{(B) Discrete inf--sup (scale inversion on the sieve)}

\begin{lemma}\label{lem:invert}
There are $c,C>0$ such that for all small $\delta$ the following holds.
Let $h=\sum_{\lambda}h_{\lambda}$ be a wavelet-packet expansion where each
$h_{\lambda}$ is a packet of Lemma~\ref{lem:scale2}-type with centre
$x_{\lambda}$ and scale $\delta_\lambda\ge\delta$ (on the sieve of \S(A),
every packet has its centre in the interior annulus, so the subfamily
below \emph{is} the whole sieve error). For the subfamily with
centres in the annulus $c_0\le\abs{x_\lambda}\le\tfrac12$ (notation
$h^{\rm bulk}=\sum_{\lambda:\,c_0\le\abs{x_\lambda}\le\frac12}h_\lambda$),
\begin{equation}\label{eq:infsup}
c\,\delta\norm{h^{\rm bulk}}_{2}
\ \le\ \norm{P_\delta Th^{\rm bulk}}_{2}
\ \le\ C\,\delta\norm{h^{\rm bulk}}_{2},
\qquad\text{and more globally}\qquad
\norm{Th}_{2}\le C\norm{h}_{2}.
\end{equation}
For origin-near modes ($\abs{x_\lambda}\le\rho\delta_\lambda$ for a fixed
$\rho$) the transport field degenerates and the gain is only
$\gtrsim\delta_\lambda^{2}$ (Lemma~\ref{lem:origin}); they are
\emph{excluded} from the sieve of \S(A) and projected into the bias, see
Remark~\ref{rem:modewise}.
\end{lemma}
\begin{lemma}[Exact radial gain near the origin]\label{lem:origin}
Let $g(x)=\varphi(\abs x/\delta)$ be radial, supported in $B_\delta$ about
the origin, $\norm{g}_2\asymp\delta^{d/2}$. Then the solution $w=Tg$ of
the radial part of \eqref{eq:DeltaT} satisfies
$w(x)=-\frac1d\int_{\abs{x}}^{\delta}s\,g(s)\,ds$ for $\abs x\le\delta$,
$w\equiv0$ for $\abs x\ge\delta$, and
$\norm{Tg}_2\asymp\delta^{2}\norm{g}_2$. More generally, for a mode
centred at distance $\abs{x_\lambda}$ with scale $\delta_\lambda$, the
mode-wise gain is at least the radial endpoint: for an origin-near
mode ($\abs{x_\lambda}\le\rho\,\delta_\lambda$, $\rho$ fixed),
\begin{equation}\label{eq:origingain}
\norm{Th_\lambda}_2\ \gtrsim\ \delta_\lambda^{2}\,\norm{h_\lambda}_2 .
\end{equation}
(used in Remark~\ref{rem:modewise} to record the degeneracy).
\end{lemma}
\begin{proof}[Proof of Lemma~\ref{lem:origin}]
For the radial bump, \eqref{eq:DeltaT} reads $\Delta w=Lg=\frac1d
\diver(xg)$, i.e.\
$\frac1{r^{d-1}}\partial_r(r^{d-1}w')=\frac1d\frac1{r^{d-1}}
\partial_r(r^{d}g)$, so $r^{d-1}w'(r)=\frac1d r^{d}g(r)+C$; regularity at
$r=0$ forces $C=0$, so $w'(r)=\frac{r}{d}g(r)$ for all $r>0$. On
$[\delta,1]$ the source $g\equiv0$, hence $w'=0$ there, and the boundary
condition $w(1)=0$ gives $w\equiv0$ on $[\delta,1]$. Continuity of $w$
at $r=\delta$ (the flux $w'$ is bounded) then yields
$w(x)=-\frac1d\int_{\abs x}^{\delta}s\,g(s)\,ds$ for $\abs x\le\delta$, and
\[
\norm{w}_2^{2}
=\abs{\mathbb S^{d-1}}\int_0^{\delta}
\Bigl|\frac1d\int_r^{\delta}sg(s)\,ds\Bigr|^{2}r^{d-1}dr
\ \asymp\ \delta^{4}\norm{g}_2^{2}
\]
by Cauchy--Schwarz (both sides $\asymp\delta^{4}\norm{g}_2^{2}$, e.g.\ the
lower bound from profiles with mass near $s=\delta$ and the upper from
$\abs{\int_r^\delta sg}| \le\delta\norm{g}_{L^2(r,\delta)}$ and the
$\delta$-length). Thus $\norm{Tg}_2\asymp\delta^{2}\norm{g}_2$.

For the one-sided bound \eqref{eq:origingain}, re-run Lemma~\ref{lem:scale2}'s
loaded-shell power count; the \emph{only} change is the shift factor
$\abs{x_m}^{2}$ in the weighted almost-orthogonality. With that proof's
notation ($\delta=\delta_\lambda$; shell
$S_k=\SET{\abs\xi\asymp\delta_\lambda^{-1}}$, volume $\asymp\delta_\lambda^{-d}$;
$\rho_k\asymp\delta_\lambda^{-1}$), the shell integrand
$\abs\xi^{-4}\abs{\widehat{Lg}}^{2}$ has its dominant
$\xi\cdot\nabla_\xi\mathbf a$ factor bounded by
$\abs{S_k}\rho_k^{2}\sum_m\abs{x_m}^{2}a_m^{2}$ (Lemma~\ref{lem:ao},
weighted part). Annulus centres have $\abs{x_m}\asymp1$; origin-near
centres have $\abs{x_m}\lesssim\delta_\lambda$, so $\abs{x_m}^{2}
\lesssim\delta_\lambda^{2}$ contributes one full extra power of
$\delta_\lambda$: the count is
\[
\norm{U}_2^{2}
\ \gtrsim\ \underbrace{\delta_\lambda^{4}}_{\abs\xi^{-4}}\,
\underbrace{\delta_\lambda^{2d}}_{\text{prefactor}}\,
\underbrace{\delta_\lambda^{-d}}_{\abs{S_k}}\,
\underbrace{\rho_k^{2}}_{\asymp\delta_\lambda^{-2}}\,
\underbrace{\delta_\lambda^{2}}_{\abs{x_m}^{2}}\sum_m a_m^{2}
\ \asymp\ \delta_\lambda^{4}\,\bigl(\delta_\lambda^{d}\sum_m a_m^{2}\bigr)
\ =\ \delta_\lambda^{4}\norm{h_\lambda}_2^{2},
\]
using $\norm{h_\lambda}_2^2\asymp\delta_\lambda^{d}\sum a_m^{2}$. The
subtracted $(\Lambda\widehat\psi)$-term lacks the $\rho_k^{2}$
amplification and is one power smaller, as in Lemma~\ref{lem:scale2}; the
harmonic correction $q=O(\delta_\lambda\norm{h_\lambda}_2)$ by
\eqref{eq:qbound} absorbs into $\delta_0$. Hence
$\norm{Th_\lambda}_2\gtrsim\delta_\lambda^{2}\norm{h_\lambda}_2$
(constants depending on $(d,\psi,c_0,\rho)$ only).
\end{proof}

\begin{proof}[Proof of Lemma~\ref{lem:invert}]
We prove the two-sided estimate for the bulk subfamily (the whole sieve
error on the interior sieve of \S(A)) and the global bound; the origin
modes it excludes are projected into $\Pi_\delta$
(Remark~\ref{rem:modewise}). Throughout, $P_\delta$ is the spectral
projection of the Dirichlet Laplacian onto
$\lambda_\nu\le C_W\,\delta^{-2}$ (spatial frequencies $\lesssim
\delta^{-1}$, matching the one-degree smoothing; see \S(A)).

\smallskip
\noindent\emph{(i) Lower bound on $\norm{Th^{\rm bulk}}_2$.} Decompose
$h^{\rm bulk}=\sum_j h_j$ into wavelet pieces at scales
$2^{-j}\ge\delta$ with centres in the annulus (each $h_j$ satisfies the
hypotheses of Lemma~\ref{lem:scale2} with scale $\delta_j=2^{-j}$ and
annulus centres). By Lemma~\ref{lem:scale2}'s lower bound applied per
scale, $\norm{Th_j}_2\gtrsim2^{-j}\norm{h_j}_2$, and the cross-terms
$\inner{Th_j}{Th_{j'}}$ for $j\ne j'$ are almost-orthogonal by
Lemma~\ref{lem:ao} (each $Th_j$ remains concentrated near scale
$2^{-j}$, both by the Fourier localization of the source and by the
Schur-summable interference off the diagonal); hence
\[
\norm{Th^{\rm bulk}}_2^{2}\ \gtrsim\ \sum_j 2^{-2j}\norm{h_j}_2^{2}
\ \ge\ \delta^{2}\sum_j\norm{h_j}_2^{2}
\ \asymp\ \delta^{2}\norm{h^{\rm bulk}}_2^{2},
\]
the last equivalence by the square-function characterization (T5). Hence
$\norm{Th^{\rm bulk}}_2\ge c\delta\norm{h^{\rm bulk}}_2$.

\smallskip
\noindent\emph{(ii) Spectral tail $\norm{(I-P_\delta)Th}_2$.} The
projector cuts at $\abs\xi\asymp\sqrt{C_W}\,\delta^{-1}$, retaining
scales $\gtrsim\delta$; the mass of $Th^{\rm bulk}$ beyond this is the
off-shell tail of Lemma~\ref{lem:scale2}. For a single scale-$\delta$
packet, \eqref{eq:Uupper}'s off-shell shells $S_j$, $2^j\gg\delta^{-1}$,
contribute a factor $2^{-N(j-k)}$ each via the Schwartz decay
$\abs{\widehat\psi(\delta\xi)}\lesssim(\delta\abs\xi)^{-N}$ (any fixed
$N$), so on $\SET{\abs\xi\gtrsim\sqrt{C_W}\delta^{-1}}$ the mass is
\[
\norm{P_{\SET{\abs\xi\gtrsim\sqrt{C_W}\delta^{-1}}}\,U}_2^{2}
\ \lesssim\ \sum_{j:\,2^{j}\gtrsim\sqrt{C_W}\delta^{-1}}2^{-N(j-k)}\,
\delta^{2}\norm{g}_2^{2}
\ \lesssim\ C_W^{-N/2}\,\delta^{2}\norm{g}_2^{2},
\]
(geometric sum over off-shell shells, $2^{j-k}\gtrsim\sqrt{C_W}$). The
harmonic correction $q$ is smooth into any Sobolev norm with constant
$O(\delta\norm g_2)$ by \eqref{eq:qbound}, so its high tail is
Schwartz-small and absorbs into the same bound. For coarser pieces
$2^{-j}\in[\delta,1]$ the cutoff $\sqrt{C_W}\delta^{-1}$ sits even
deeper in the off-shell tail, so the bound is uniform; summing over
scales as in step~(i),
\begin{equation}\label{eq:spectail}
\norm{(I-P_\delta)Th^{\rm bulk}}_2
\ \le\ \widetilde C\,C_W^{-N/4}\,\delta\norm{h^{\rm bulk}}_2,
\end{equation}
with $\widetilde C=\widetilde C(d,\psi,c_0)$. Choosing the fixed constant
$C_W$ (hence the sieve $V_\delta$) so large that $\widetilde
C\,C_W^{-N/4}\le c/4$ ($N\ge4$ fixed, $C_W$ independent of $\delta$)
makes the tail $\le\tfrac{c}{4}\delta\norm{h^{\rm bulk}}_2$.

\smallskip
\noindent\emph{(iii) Combination.} By the triangle inequality (ordinary
norm subtraction),
\[
\norm{P_\delta Th^{\rm bulk}}_2
\ \ge\ \norm{Th^{\rm bulk}}_2-\norm{(I-P_\delta)Th^{\rm bulk}}_2
\ \ge\ c\,\delta\norm{h^{\rm bulk}}_2-\tfrac c4\,\delta
\norm{h^{\rm bulk}}_2
\ =\ \tfrac{3c}{4}\,\delta\norm{h^{\rm bulk}}_2 ,
\]
combining step~(i) and \eqref{eq:spectail}; both terms are
$O(\delta\norm{\cdot}_2)$ and the tail constant $c/4<c$, so the
difference is positive. The upper bound
$\norm{P_\delta Th^{\rm bulk}}_2\le\norm{Th^{\rm bulk}}_2
\lesssim\delta\norm{h^{\rm bulk}}_2$ and the global
$\norm{Th}_2\le C\norm{h}_2$ follow from Lemma~\ref{lem:scale2} /
Remark~\ref{rem:degree1}.
\end{proof}

\begin{remark}[Origin content and the interior sieve]\label{rem:modewise}
Lemma~\ref{lem:origin} records that the transport field's degeneracy makes
the origin the \emph{least} revealing region (gain only
$\asymp\delta_\lambda^{2}$ there, versus $\asymp\delta_\lambda$ in the bulk
annulus). Since an adversarial $h$ could concentrate at the origin ---
where inversion against the $O(\delta)$ noise/bias terms would lose a
factor $\delta^{-1}$ --- the sieve $\F_\delta$ is restricted to
bulk-annulus packets, and the origin/boundary content is projected into
$\Pi_\delta f_0$ (hence into the bias, whose $u$-image is handled by the
degree-one gain of \eqref{eq:finebias}). Accordingly, in \S(C) the lower
bound of \eqref{eq:infsup} is applied to the \emph{whole} sieve error $h$,
supported in the bulk annulus by construction, and the analysis uses the
annulus two-sided estimate of Lemma~\ref{lem:scale2}, valid for all
$\alpha>1$, $d\ge2$ (Lemmas~\ref{lem:scale2}; $c_0,c_\psi,C_E''$
absorbed).
\end{remark}

\subsection{(C) Risk bound}
Let $\delta=\delta_*$ at the end; for now $\delta>0$ general. Decompose,
with the comparison point $m:=\Pi_\delta f_0$, the bias $e=m-f_0$, and the
sieve error $h:=\widehat f-m\in\F_\delta-\F_\delta$,
\[
\widehat f-f_0=h+e .
\]
\textbf{Bias.} By \eqref{eq:approx}, $\norm{e}_2^{2}\lesssim\delta^{2\alpha}$.

\noindent\textbf{Linearize about the comparison point (multiplicative
$r$).} Rather than linearize about $\one$ (which produces an
\emph{additive} remainder $\sim r\norm{h}_2$ that cannot be absorbed at
fixed $r$), we linearize about $m=\Pi_\delta f_0$. Note
$1-2r\le m\le1+2r$ and $\norm{\nabla u_m}_{\infty}\lesssim1$ ((T6), since
$m\in\F(\alpha,2r)\subset C^\alpha$). Write
$u_{\widehat f}-u_{m}=T_{m}h+\widetilde\rho$, where $T_{m}$ is the
linearized forward map at $m$ and $\widetilde\rho$ is its quadratic
remainder. The same energy argument as Lemma~\ref{lem:quad}, with
$\nabla u_m$ in place of $\nabla u_0$, gives
$\norm{\widetilde\rho}_2\lesssim\norm{h}_\infty\norm{h}_2$. The
linearization about $\one$, $T=T_\one$, differs from $T_m$ by an
operator-difference reflecting the $r$-closeness of $m$ to $\one$; the
essential point is that this enters \emph{multiplicatively}: for
wavelet-packet $h$ on the sieve at scale $\delta$ the $h$-dependent part
of $T_m-T$ satisfies
\begin{equation}\label{eq:opdiff}
T_mh-Th=-\Delta^{-1}\diver\bigl(h\nabla(u_m-u_0)\bigr),
\qquad
\norm{T_mh-Th}_2\ \lesssim\ r\,\delta\,\norm{h}_2 ,
\qquad m-\one=O(r)\text{ in }C^\alpha,
\end{equation}
To see the identity in \eqref{eq:opdiff}, subtract the linearized
equations: $\diver(\nabla T_mh)=-\diver(h\nabla u_m)$ and
$\diver(\nabla Th)=-\diver(h\nabla u_0)$ give
$\diver\nabla(T_mh-Th)=-\diver\bigl(h\nabla(u_m-u_0)\bigr)$, i.e.\
$T_mh-Th=-\Delta^{-1}\diver\bigl(h\nabla(u_m-u_0)\bigr)$ (this is a
first-order expansion of the \emph{difference} $(T_m-T)h$, the
$h$-independent piece $(T_m-T)0=u_m-u_0$ being exactly accounted for by
$u_m-u_0=\Delta^{-1}\diver((m-\one)\nabla u_0)$).
For the bound, the only subtlety is the boundary layer of $m-\one$: by the
interior Schauder estimate
$\norm{u_m-u_0}_{C^{2,\beta}(K)}\le C_K\norm{m-\one}_{C^\alpha}$ on
compact $K\Subset D$ (the right side of
$\diver(m\nabla(u_m-u_0))=\diver((m-\one)\nabla u_0)$ is the divergence
of the vector field $(m-\one)x/d$, whose $C^{\alpha-1}(K)$ norm is
$\lesssim_K r$; (T6) applies) and the packet $h$ is supported in
$\SET{\abs x\le\tfrac34}=:K$, so
$\norm{h\nabla(u_m-u_0)}_2\le C_K r\norm h_2$.
Repeating the
single-scale Fourier computation of Lemma~\ref{lem:scale2} with the
vector source $h\nabla(u_m-u_0)$ in place of $xg$ --- for the
\emph{upper} bound only, where the estimate reduces to
$\norm{-\Delta^{-1}\diver F}_2\lesssim\delta\norm F_2$ for any vector
field $F$ whose spatial scale is $\delta$ (Lemma~\ref{lem:scale2}'s
upper bound \eqref{eq:Uupper} applied componentwise; the slow factor
$\nabla(u_m-u_0)$ is $O(r)$ in $L^\infty$ and $C^{\alpha-1}$ and does not
change the scale) --- yields
\[
\norm{T_mh-Th}_2
\ \lesssim\ \delta\,C_r\,\norm{h}_2\ \lesssim\ r\,\delta\,\norm{h}_2
\]
(this is the multiplicative $r$: the estimate scales linearly with
$\norm{m-\one}_{C^\alpha}\lesssim r$, and fails for $r=1$). The remaining
piece $(u_m-u_0)$ of $(T_m-T)$ is \emph{independent of $h$} and is
grouped, below, with the bias/discrepancy term (it acts as an additional
$O(r)$ offset inside $u_m-u_{f_0}$; its contribution to the risk is
handled with the bias in the accounting display below).
Consequently the signal satisfies the recoverable lower bound
\begin{equation}\label{eq:infsupapp}
\norm{P_\delta T_mh}_2
\ \ge\ \norm{P_\delta Th}_2-\norm{(T_m-T)h}_2
\ \ge\ \bigl(c-Cr\bigr)\delta\norm h_2
\ \ge\ \tfrac c2\,\delta\norm h_2
\qquad\text{for }r\le r_0:=\tfrac{c}{2C},
\end{equation}
by Lemma~\ref{lem:invert} applied to the annulus packet $Th$ and
\eqref{eq:opdiff}. Since the sieve perturbations $h$ are by
construction supported in the bulk annulus (no origin-near modes are
estimated; Lemma~\ref{lem:invert}'s lower bound applies to the
\emph{whole} $h$), this is the recovery lower bound used in
\eqref{eq:basic}; on the origin content absorbed into $\Pi_\delta$ see
Remark~\ref{rem:modewise}.

\par\noindent\textbf{Absorption via per-scale Nikolskii on the compact sieve.}
On the compact sieve $\F_\delta\subset\SET{\norm{\cdot}_{C^\alpha}\le2r}$,
a difference $h=\sum_j h_j$ of two sieve elements has
$\norm{h_j}_{C^\alpha}\le4r$ at each scale $2^{-j}\ge\delta$. The
per-scale Nikolskii inequality and the $C^\alpha$ class cap give
$\norm{h_j}_\infty\lesssim\min\bigl(2^{jd/2}\norm{h_j}_2,\,
2^{-j\alpha}r\bigr)$. Summing over the $\lesssim\log(1/\delta)$ scales
$2^{-j}\ge\delta$, and using Lemma~\ref{lem:quad}'s
$\norm{\widetilde\rho}_2\lesssim\norm{h}_\infty\norm{h}_2$ together with
the pointwise triangle inequality $\norm{h}_\infty\le\sum_j\norm{h_j}_\infty$
(the pieces are summable in $L^\infty$ since $\norm{h_j}_\infty\lesssim
2^{-j\alpha}r$ and $\alpha>0$),
\[
\frac{\norm{\widetilde\rho}_2}{\delta\norm{h}_2}
\ \le\ \frac{\norm{h}_\infty}{\delta}
\ \le\ \delta^{-1}\sum_{2^{-j}\ge\delta}\norm{h_j}_\infty
\ \lesssim\ \delta^{-1}\sum_{2^{-j}\ge\delta}
\min\bigl(2^{jd/2}\norm{h_j}_2,\,2^{-j\alpha}r\bigr).
\]
Using the bootstrap $\norm{h}_2\le2r$ and the uniform $C^\alpha$ bound,
each summand is $\lesssim2^{-j\alpha}r$, so
$\delta^{-1}\sum_{j}2^{-j\alpha}r\lesssim\delta^{-1}r\delta^{\alpha}
=r\delta^{\alpha-1}\to0$ \emph{because} $\alpha>1$. Hence, at the scale of
interest and quantitatively (no undefined exponent),
\begin{equation}\label{eq:absorb}
\norm{\widetilde\rho}_2\ \le\ r\,\delta^{\alpha-1}\bigl(\delta\norm{h}_2\bigr)
\ =:\ \omega_\delta\,\delta\norm{h}_2,
\qquad \omega_\delta:=r\delta^{\alpha-1}\to0\ \ (\alpha>1).
\end{equation}
This is the genuine multiplicative absorption, valid for all $\alpha>1$,
all $d\ge2$.

\par\noindent\textbf{Basic inequality and risk.} Write, abusing notation,
$\Upsilon(g)=\norm{P_\delta u_g-Y_\delta}_2^{2}$,
$\Delta Y:=P_\delta(u_{f_0}-Y_\delta)=n^{-1/2}\xi_\delta$ (the projected
negative noise), and set
\[
X\ :=\ P_\delta\bigl(u_{\widehat f}-u_m\bigr),\qquad
Z\ :=\ P_\delta\bigl(u_{f_0}-u_m\bigr).
\]
From $\Upsilon(\widehat f)\le\Upsilon(m)$, expanding the squares as
$\norm{X-(\Delta Y+Z)}_2^{2}\le\norm{\Delta Y+Z}_2^{2}$ and cancelling
the common $\norm{\Delta Y+Z}_2^{2}$,
\[
\norm{X}_2^{2}\ \le\ 2\,\bigl|\inner{X}{\Delta Y+Z}\bigr|
\ \le\ 2\norm X_2\bigl(\norm{\Delta Y}_2+\norm Z_2\bigr),
\]
so (dividing by $\norm X_2$; if $X=0$ the conclusion is trivial)
\begin{equation}\label{eq:basic0}
\norm{X}_2\ \le\ 2\,\norm{\Delta Y}_2+2\,\norm Z_2
\ =\ 2n^{-1/2}\norm{\xi_\delta}_2+2\norm{P_\delta(u_{f_0}-u_m)}_2 .
\end{equation}
On the left, $u_{\widehat f}-u_m=T_mh+\widetilde\rho$ with
$\norm{\widetilde\rho}_2\le\omega_\delta\,\delta\norm h_2$ by
\eqref{eq:absorb}, hence by \eqref{eq:infsupapp}
\[
\norm{X}_2\ \ge\ \norm{P_\delta T_mh}_2-\norm{\widetilde\rho}_2
\ \ge\ \tfrac c2\delta\norm h_2-\omega_\delta\delta\norm h_2
\ \ge\ \tfrac c4\,\delta\norm h_2
\qquad(\delta\le\delta_0).
\]
For the right side of \eqref{eq:basic0}, the \emph{discrepancy}
$u_{f_0}-u_m$ is the solution error generated by the approximation
residual $f_0-m=f_0-\Pi_\delta f_0=-e$; by the fine-scale bias bound
(next paragraph),
$\norm{u_{f_0}-u_m}_2\lesssim\delta^{1+\alpha}$, i.e.\
$\norm Z_2=\norm{P_\delta(u_{f_0}-u_m)}_2
\lesssim\delta^{1+\alpha}$.
Combining the last two displays with \eqref{eq:basic0},
\begin{equation}\label{eq:basic}
\tfrac{c}{4}\delta\norm{h}_2
\ \lesssim\ n^{-1/2}\norm{\xi_\delta}_2+\delta^{1+\alpha}
+\omega_\delta\,\delta\norm{h}_2 ,
\end{equation}
and the $\omega_\delta\delta\norm{h}_2$ term is absorbed on the left for
$\delta\le\delta_0$ (since $\omega_\delta\to0$).

\par\noindent\textbf{Fine-scale bias.} Write $w:=u_{f_0}-u_m$; both
$u_{f_0}$ and $u_m$ vanish on $\partial D$ (the Dirichlet condition of
the very-weak formulation \eqref{eq:veryweak}), so $w\in H^{1}_0$ solves
$\diver(f_0\nabla w)=\diver G$, $G:=e\nabla u_m$, with
$\norm{G}_2\le\norm e_2\norm{\nabla u_m}_\infty\lesssim\delta^\alpha$
by \eqref{eq:approx} and $\norm{\nabla u_m}_\infty\lesssim1$ (T6).
Split $G=G_{\rm int}+G_{\rm col}$, $G_{\rm col}=G\,\mathbf 1_{S}$ the
part in a collar $S=\SET{\abs x>1-s}$ of width $s$ around $\partial
D$. Two points:
(i) \emph{collar remainder}: by H\"older and \eqref{eq:approx},
$\norm{e}_{L^{2}(S)}\le\norm{e}_{L^{p}(S)}\abs{S}^{\frac12-\frac1p}
\lesssim s^{\frac12-\frac1p}\delta^{\alpha}$, with $p>2$ taken large
(the $C^\alpha$-cap tail is Nikolskii-small,
$\norm{e}_{L^{p}(S)}\lesssim\delta^{\alpha}$, the projected region
$O(r)$); for $s\lesssim\delta$ the collar is
$\lesssim\delta^{\alpha+\frac14}=o(\delta^{\alpha})$, and no
$r\delta^\alpha$ term appears --- collar content is projected into
$\Pi_\delta$;
(ii) \emph{interior gain}: each scale-$\delta$ bump of $e$, at
distance $\gtrsim s$ from $\partial D$, is mean-zero, so by the exact
Fourier mechanism of Lemma~\ref{lem:scale2} (the
$\abs\xi^{-1}$-gain on the loaded shell; the harmonic correction
\eqref{eq:qbound} applies since $\supp G_{\rm int}$ is interior)
$\norm{G_{\rm int}}_{H^{-1}}\lesssim\delta\norm e_2\lesssim
\delta^{\alpha+1}$. Combining, testing
against $\phi=-\Delta^{-1}_{f_0}\psi$ as in Lemma~\ref{lem:apriori}
and using Poincar\'e on $w\in H^{1}_0$,
\begin{equation}\label{eq:finebias}
\norm{w}_2=\norm{u_{f_0}-u_m}_2
\ \lesssim\ \norm{G}_{H^{-1}}
\ \lesssim\ \delta\norm{e}_2+\delta^{\alpha+\frac14}
\ \lesssim\ \delta^{1+\alpha}.
\end{equation}
Thus
$\norm Z_2=\norm{P_\delta(u_{f_0}-u_m)}_2
\lesssim\delta^{1+\alpha}$.

\par\noindent\textbf{Noise enters linearly (no union bound necessary).} In the
comparison inequality the noise appears only through the \emph{linear}
functional $h\mapsto n^{-1/2}\inner{\xi_\delta}{P_\delta Th}$ evaluated at
the random point $h=\widehat f-\Pi_\delta f_0$; by Cauchy--Schwarz,
$n^{-1/2}\abs{\inner{\xi_\delta}{P_\delta Th}}\le
n^{-1/2}\norm{\xi_\delta}_2\,\norm{P_\delta Th}_2$, and after division by
$\norm{P_\delta Th}_2$ this contributes $n^{-1/2}\norm{\xi_\delta}_2$ on
the event $\SET{\widehat f\in\F_\delta}$, uniformly. Because the projected
model is finite-dimensional and the functional is linear,
$\E\norm{\xi_\delta}_2^2=\dim V_\delta\asymp\delta^{-d}$, and \emph{no
supremum / union bound over the sieve is needed here}: a logarithmic factor
would arise only if one insisted on a \emph{uniform}-in-$g$ comparison,
which the least-squares basic inequality avoids. Account for every term
of \eqref{eq:basic}: after absorbing $\omega_\delta\delta\norm h_2$ and
dividing by $\delta$, $\norm h_2\lesssim\delta^{-1}\bigl(n^{-1/2}
\norm{\xi_\delta}_2+\delta^{1+\alpha}\bigr)$, so squaring (and using
$\E\norm{\xi_\delta}_2^2\asymp\delta^{-d}$)
$\E\norm h_2^2\lesssim\delta^{-d-2}/n+\delta^{2\alpha}$: the
discrepancy $\delta^{1+\alpha}$ contributes $\delta^{2\alpha}$, not
$r^{2}\delta^{2\alpha-2}$ --- the boundary collar/collar content
responsible for the latter has been projected into the bias, whose
$u$-image uses the interior degree-one gain \eqref{eq:finebias}. With
the balance
$\delta=\delta_*=n^{-1/(2\alpha+d+2)}$ chosen so
$\delta_*^{-d-2}/n=\delta_*^{2\alpha}$, one therefore obtains the
\emph{log-free} estimate
\[
\E\norm{\widehat f-f_0}_2^{2}
\lesssim\norm{e}_2^{2}+\E\norm{h}_2^{2}
\lesssim\delta_*^{2\alpha}+\frac{\delta_*^{-d-2}}{n}
\asymp n^{-\frac{2\alpha}{2\alpha+d+2}} .
\]
All contributing terms (bias, discrepancy, noise) are
$\lesssim\delta_*^{2\alpha}$ at the balance; the upper bound is free of
logarithmic loss.

\section{Conclusion and consistency checks}

\begin{itemize}[itemsep=2pt]
\item \textbf{Rate.} Lower bound (Proposition~\ref{prop:lower}):
$n^{-2\alpha/(2\alpha+d+2)}$ (no log). Upper bound
(Proposition~\ref{prop:upper}): $n^{-2\alpha/(2\alpha+d+2)}$ (no log:
the noise enters the least-squares basic inequality only through a
\emph{linear} functional, so no union bound is needed). The two agree, so
the minimax rate is $n^{-2\alpha/(2\alpha+d+2)}$ up to logarithmic factors.
\item \textbf{Interpretation.} The exponent denominator
$2\alpha+d+2$ exceeds the direct-problem $2\alpha+d$ by exactly $2$,
reflecting \emph{one} degree of smoothing of the linearized inverse $T$
(order $-1$), a consequence of the transport structure
$\Delta T=I+\frac{x\cdot\nabla}{d}$ and the $L^{2}$-invertibility of
$I+\frac{x\cdot\nabla}{d}$ on $L^{2}(D)$
(Lemmas~\ref{lem:lin},~\ref{lem:Linv} and the exact Fourier identity
Lemma~\ref{lem:fourier}).
\item \textbf{Role of $\alpha>1$.} Used decisively in both bounds to make
the quadratic remainder $R(h)$
(Lemma~\ref{lem:quad}) negligible against the linear signal $Th$ at the
minimax scale (ratio $\asymp\delta_*^{\alpha-1}\to0$), cf.\
Remark~\ref{rem:alpha>1}. For $\alpha\le1$ the linearization is not
faithful at scale $\delta_*$; the rate would then require a genuinely
nonlinear analysis (we make no claim in that regime).
\item \textbf{Role of small $r$.} $r\le r_0(\alpha,d)$ is used for
(i) uniform ellipticity $1-r\le f\le1+r$, (ii) absorbing the quadratic
remainder into the linear signal, and (iii) the uniform apriori bound
$\norm{u_f}_\infty\lesssim1$. All structural constants and $r_0$ depend only on $(\alpha,d,\psi)$;
$C$ is uniform in $r\in(0,r_0]$ and $n$; the lower-bound $c$ carries the
$r^{2}$ factor of the shrinking class.
\item \textbf{Model handling.} $u_f$ is defined by the very-weak formulation
\eqref{eq:veryweak}; no $H^1$ membership is asserted
(Section~\ref{sec:forward}). The white noise is only ever projected onto
the finite-dimensional sieve $V_\delta$, where it is a well-defined
Gaussian vector, so the estimator/risk are rigorous on the white-noise
model.
\end{itemize}

\end{document}
